% ============================================================
%  R. Nielsen, "Om det oprindelige Forhold mellem Religion og Videnskab"
%  Kjøbenhavn: J. H. Schultz, 1881
%  Indbydelsesskrift til Kjøbenhavns Universitets Fest, 8. april 1881
%  TRANSCRIPTION (Danish) — from Royal Danish Library digital scan
% ============================================================
%
%  EDITION RECORD (2026-09-27)
%  Scan: KB digitisation, 83 PDF pp., sha256 8f9038aa...c3ed. The essay is
%  printed pp. 1-55 = PDF 12-66; PDF 67 onward (the university's prize
%  reports) is not Nielsen's and is not transcribed.
%  Page map: PDF = printed + 11, uniform, verified by eye on every folio.
%  \opage{N} stands at the word where printed page N begins; pp. 2-8 were
%  first placed by machine alignment and have been checked at the image.
%  Conventions: diplomatic; 1881 orthography as printed; quotation marks
%  „...“ as printed; letterspacing -> \emph; Latin, German and French are
%  set in the same roman as the Danish, so no \textit; Greek in Unicode
%  (textalpha); footnotes are marked *) and restart on each printed page,
%  set here as \footnote with the printed mark in a comment; the section
%  numerals I.-IV. are centred as printed.
%  Printer's defects are transcribed as printed and logged in a "% sic"
%  comment at the site. The print's quotation marks do not pair on pp.
%  20, 38, 39, 46 and 48 (logged there); the counts happen to net to zero.
%  Word-accuracy: pp. 8-55 were transcribed in 4 batches, every batch
%  diffed against the scan's embedded text layer before splicing
%  (ocrdiff.py): 15 candidates, 0 transcription errors, 0 unresolved.
%  pp. 1-8, transcribed earlier without a word check, were collated
%  against the page images on 2026-09-27: 17 word errors and 11 other
%  faults (punctuation, emphasis, paragraphing, one page marker)
%  corrected. See TRANSCRIPTION-PLAYBOOK.md §3 step 4.
% ============================================================
\documentclass[12pt,a4paper]{article}
\usepackage[utf8]{inputenc}
\usepackage[T1]{fontenc}
\usepackage{libertinus}
\usepackage{libertinust1math}
\usepackage[danish]{babel}
\usepackage{textalpha}
\usepackage[protrusion=true,expansion=true]{microtype}
\usepackage{setspace}
\usepackage[margin=1.2in]{geometry}
\usepackage{fancyhdr}
\usepackage{hyperref}

\hypersetup{
  pdftitle={Om det oprindelige Forhold mellem Religion og Videnskab},
  pdfauthor={R. Nielsen},
  hidelinks
}

\pagestyle{fancy}
\fancyhf{}
\rhead{Nielsen, \emph{Om det oprindelige Forhold} (1881)}
\cfoot{\thepage}

\setstretch{1.3}

\title{\textbf{Om det oprindelige Forhold\\[4pt]
       mellem Religion og Videnskab}}
\author{R.\ Nielsen}
\date{Kjøbenhavn: J.\ H.\ Schultz, 1881}

% ---- original-page markers (paginate.py) ----
\usepackage{marginnote}
\newcommand{\opage}[1]{\marginnote{\footnotesize\textit{[#1]}}}
\newcommand{\apage}[1]{\marginnote{\footnotesize\textit{[#1]}}}
\begin{document}
\maketitle
\thispagestyle{fancy}

\bigskip

\begin{center}I.\end{center}

\noindent Den bittre Strid mellem Religion og Videnskab vil ikke lade
sig hæve, med mindre det skeer paa den ufravigelige Betingelse,
at begges oprindelige Selvstændighed gjensidig anerkjendes. Religionen
for sig maa da kunne bestaae uden al Videnskab, og Videnskaben for sig
uden al Religion; men hvorledes er det muligt? At Religionen ikke er
af videnskabelig Oprindelse, „den aabenbarede Religion“ saa lidt som
Mytherne, maa vel indrømmes; men Historien lærer, at Videnskaben neppe
slaaer Øinene op, før den begynder at rokke ved den givne Religion og
indlede Striden. Hos Grækerne var det især Philosopherne, som førte
Angrebet. Xenophanes paaviste Modsigelser i den polytheistiske
Folketro; imod de mange Guder satte han Eenheden, imod deres timelige
Oprindelse Evigheden, mod det Ubestandige i deres Væsen
Uforanderligheden, imod deres Lighed med Mennesker Ophøietheden
o.\ s.\ v.\footnote{Zeller, Die Philosophie der Griechen, 1ster
Theil, 2te Auflage, Tübingen 1856. S.\ 381.}. % *)
Hos de store Digtere
sporer man en trinviis Forandring i Gudernes Væsen; hos Euripides begynde
de allerede at hjemfalde til Reflexionen. Den naturlige Forstaaelse
rykker den Dannede nærmere. Kunsten holdt paa Formen, men netop den
plastiske Formfuldendelse tærede paa Gudernes Realitet. Zeus i Olympia
stod vistnok som Mægler mellem Folketroen og den philosophisk-kritiske
Opfattelse; men Indtrykket var æsthetisk, og det overvældende æsthetiske
Indtryk kvalte det Religiøse. Man \opage{2} kan uden Overdrivelse paastaae, at
den videnskabelige Tænkning, som undergravede de græske Guders
Autoritet, havde en mægtig Forbundsfælle i den plastiske Kunst.

Hvad Forholdet mellem det Gamle og Nye Testamentes aabenbarede Religion
angaaer, stiller Sagen sig anderledes. Begavet med rige Evner, dyb
Følelse, storartet Phantasi, skarp Forstand og stærke Lidenskaber, var
Israels Folk dog ikke som Grækerne særlig anlagt enten paa plastisk
Kunst eller dybtgaaende videnskabelig Forskning; det var i
Verdenshistorien Jehovas Folk, i særegen Betydning Religionens Folk,
Aabenbaringens udvalgte Folk. En kunstnerisk Efterligning af det
Guddommelige var forbudt, Poesien stod udelukkende i Religionens
Tjeneste, og den kritisk prøvende Videnskab havde ingen Rod i
Folkenaturen. At Moseloven med sin høiere Ethik ikke tilsteder nogen
grundig videnskabelig Undersøgelse, er indlysende. Hvor meget de ti
Bud end synes at tiltale den menneskelige Fornuft, saa er det dog strax
iøinefaldende, at det ikke er Fornuftens Grunde, men Jehovas Ord,
hvorpaa Eftertrykket falder. Man lægge vel Mærke til, at den samme
Gud, der aabenbarer sin Villie i den almindelige Lov, ogsaa udtaler sig
directe i de særlige Love.
Angaaende Tabernaklet siger Jehova (2den Moseb.\ XXV): „Og dette er
Offergaven, som I skulle tage af dem: Guld og Sølv og Kobber; Mørkeblaat
og Purpurrødt og Carmosinrødt og fiint Linned og Gedehaar“ o.\ s.\ v.
„Og Tabernaklet (XXVI) skal du gjøre af ti Forhæng \ldots\ Længden af
hvert Forhæng skal være otte og tyve Alen og Breden fire Alen paa hvert
Forhæng, eet Maal for alle Forhæng“ \ldots\ Ingen skal sige, at disse
specielle Bud overskride Fornuftens Fatteevne; det Overnaturlige, for
os Ufattelige ligger tvertimod deri, at den Almægtige, Himlens og
Jordens Skaber, har talt saa nøiagtigt om Tepper og Sløifer og heri
slet Intet overladt til Moses' sunde Sands. Det er dette, som her maa
forekomme Fornuften aldeles ubegribeligt. Paa modsat Maade ydmyges
Fornuften ikke blot i det Gamle, men ogsaa i det Nye Testamente, naar
det Overnaturlige slaaer igjennem saa pludselig, saa vældig, at Naturen
revner, Lovenes Baand brister og Miraklet \opage{3} skeer. Eller staaer maaskee
Begivenheden med Gadarenernes Sviin ikke lige saa høit over Fornuften,
som Miraklet (2den Kong.\ XIII) med den døde Mand, der faldt ned i
Propheten Elisas Grav: „der kom Liv i ham, og han reiste sig op paa
sine Been“. Der viser sig i slige Træk en saadan Splid mellem den
guddommelige Aabenbaring og den menneskelige Fornuft, at ethvert
Forlig mellem Religion og Videnskab maa synes umuligt. Skal
Aabenbaringen gjælde, saa maa Videnskaben bort; skal Videnskaben
gjælde, saa maa Aabenbaringen opløses: et fortvivlet Dilemma.

Vi ville her ikke opholde os ved enkelte Træk af de Stridigheder, som
nu i næsten to tusinde Aar ere førte mellem Aabenbaringstroen og
Fornuften, Stridigheder, hvorunder snart Religionen har løftet det
kirkelige Banner høit og været i Færd med at underkue Videnskaben,
snart igjen den seirende Videnskab paa den krænkede Fornufts Vegne
været i Begreb med at opløse Religionen; vor Opgave er at gaae tilbage
til Stridens oprindelige Grund, for at undersøge, om der i Grunden
selv ikke skulde skjule sig en Misforstaaelse, en uklar Sammenblanding
og Forvexling af religiøs og videnskabelig Erkjendelse. Christi Ord
(Joh.\ VIII, 32): „\emph{I skulle erkjende Sandheden og Sandheden skal
frigjøre Eder}“, hæve sig høit over alle Partistridigheder, og maae da
lige saa fuldt komme Videnskaben som Religionen til gode.

Men, indvender man, den Sandhed, der skal frigjøre os, maa nødvendig
stemme overeens med sig selv i Et og Alt. Hvis Noget, seet fra en
Side, vel er sandt, fra en anden Side derimod usandt: da ligger dette
ikke i Sandheden, men i dens Fremstilling og Opfattelse. Er end al vor
Erkjendelse Stykværk, saa kan det til Sandheden selv svarende
Erkjendelsesprincip dog umulig lade sig udstykke. To indbyrdes
stridende Erkjendelsesprinciper maae jo forudsætte to principielt
modsatte, mod hinanden stridende Sandheder, men det er aldeles
urimeligt. Dersom den religiøse Erkjendelse staaer i aabenbar Strid
med den videnskabelige, da kan dette ikke betyde, at her virkelig er
to Sandheder og for Erkjendelsen to absolut ueensartede Principer; men
et af de to \opage{4} foregivne Principer maa nødvendig være falsk, enten det
religiøse eller det videnskabelige. Har det religiøse Princip Gyldighed,
saa maa det omfatte hele Sandheden, ogsaa den videnskabelige; Videnskaben
kan da ikke gjøre Fordring paa at have noget særegent Princip. Er det
omvendt det videnskabelige Princip, der har Gyldighed, da er Religionen
i sig selv principløs, og maa i Et og Alt være Videnskabens Dom
underlagt. --- Vi skulle i det Følgende undersøge, hvorledes det hermed
forholder sig.

At Sandheden i evig Forstand maa være een og udeelbar, ville vistnok
Alle indrømme; men Sandhedens Væsen er i sig selv dobbelsidig,
hverken blot noget Objectivt ei heller noget blot Subjectivt, men en
Erkjendelsens Eenhed af begge. Uden erkjendende Subject gives der ingen
Erkjendelse, og uden et Object, som bliver erkjendt, gives der heller
ingen Erkjendelse. Sandheden selv kan altsaa hverken være det
Subjective for sig ei heller det Objective for sig; den gjennemtrængende
Eenhed af Subjectivt og Objectivt maa gaae igjennem Altilværelsen,
omfattende den hele Virkelighed med alle dens Muligheder i det Mindste
som i det Største. Dette maa Mennesket med sin ufuldkomne Erkjendelse
indrømme; men hermed indrømmer det tillige, at det i sin Begrændsning
umulig kan gjennemskue den absolute Sandheds Indhold i dettes utallige
Former og Concretioner. Skal Mennesket virkelig erkjende Sandheden, da
maa Sandheden dele sig; og den deler sig som Lyset. Idet Sandhedens
Lys straaler ind i Menneskets Bevidsthed, afsætter det to
Erkjendelses-Midtpunkter, to Principer, som hvert for sig antage sit
relative Præg af det Absolute, nemlig Troens og Videns; det første
angaaer væsenlig Menneskets eget Selv, dets Indre, og er personligt;
det andet angaaer Tingene og deres væsenlige Sammenhæng, og er for saa
vidt upersonligt.

Skiller man Tro og Viden fra de tilsvarende Principer, ville de
bestandig flyde over i hinanden. At Troen umulig kan undvære al Viden,
er lige saa indlysende som at Viden ikke kan undvære al Tro. Den
Troende maa jo nødvendig vide, hvad det er, hvorpaa han troer og hvad
det at troe vil sige; lige saa maa den Vidende troe paa sine Sandser og
sin Forstand. Skeptikeren, der \opage{5} hverken kan troe paa det Ene eller det
Andet, ei en Gang paa sin egen Skepticisme, efterdi Sandheden er, at
der ingen Sandhed er, er i lige Grad afspærret fra Videnskaben og fra
Religionen.

Der er en dybt gaaende Forskjel mellem religiøs og videnskabelig Tro,
mellem religiøs og videnskabelig Viden. Den religiøse Tro har et
psychisk-ethisk Præg og er væsenlig en Samvittighedssag, et reent
personligt Anliggende; den videnskabelige Tro angaaer ikke Mennesket
særligt, men udtrykker Tillid til Videnskabens Grundbetingelser og
Realitet. Den religiøse Tro sætter Kundskab om Gud og det Overnaturlige
langt høiere end Kjendskab til Virkelighedens Love og de naturlige
Ting: al tilsvarende Viden faaer herved alene sin Betydning. Det
Overnaturlige kan ikke sandses, og maa dog give sig tilkjende paa
umiddelbar Maade; det antager da en sandselig-oversandselig Form i
Phantasien, og gjør derigjennem Indtryk paa Følelserne: det aabenbarer
sig. Athene aabenbarede sig for Achileus (Il.\ I, 194--98), Jehova
aabenbarer sig for Abraham, for Moses, for Propheterne; selv Guds
Aabenbaring i Christus gjør herfra ingen Undtagelse. I Johannes første
Brev I, 1 hedder det vel: „Det, som var fra Begyndelsen, det, som vi
have hørt, det, som vi have seet med vore Øine, det, som vi have
beskuet, og vore Hænder have følt paa, nemlig om Livets Ord“; det
kunde vel synes, som blev det Overnaturlige her ligefrem bekræftet af
Sandserne og bevidnet paa sandselig Maade; men sammenligne vi hermed
(Matth.\ XVI, 15--17), hvor Jesus siger til Peder, der erkjender ham
for Guds Søn, „salig er du Simon, Jonas's Søn, thi det har Kjød og
Blod ikke aabenbaret dig, men min Fader, som er i Himlen“, saa komme
vi til fuld Forvisning om, at en sand Viden om det Overnaturlige maa
være af overnaturlig Oprindelse. Den Viden, der har den religiøse Tro
til Grundlag, kan altsaa ikke opløse Troen; naar Troen svinder, blive
de overnaturlige Objecter øieblikkelig forvandlede til Myther og
Indbildninger, Jehova og Christus lige saa vel som Zeus og Athene. I
denne Henseende har Anselmus med sit credam, ut intelligam % sic: the print has „credam“ (Anselm's formula is credo ut intelligam), p. 5
unegtelig Ret imod Abailard, der vilde vende Forholdet om.

\opage{6} Anderledes i Videnskaben. Vel kan man ogsaa sige, at den videnskabelige
Viden begynder med Tro; men deels er Visheden om det for Bevidstheden
Nærværende ikke Tro, men en umiddelbar Viden; deels er Alt hvad der
foreløbig antages i Tro, fordi det antages i uprøvet Tillid, bestemt
til at underkastes en Prøvelse, hvorved Troen, ifølge det Cartesianske:
de omnibus dubitandum est, opløser sig i Tvivl, og naar saa
Tvivlen er hævet ved Begrundelsen, forvandler sig til Viden. Vi begynde
med at orientere os i Verden; Sandser, Forstand og Fornuft virke
sammen paa umiddelbar Maade; snart er det Sandserne, snart Forstanden,
der har eensidig Overvægt, medens Fornuften dog altid kræver, at
Vexelbestemmelsen af det, som sandses, og det, som tænkes, holdes i
Hævd, thi herpaa beroer den videnskabelige Viden. Idet Undersøgelsen
begynder angaaende Forskjellen mellem hvad der i Tilværelsen er
Foranderligt og hvad der som et Uforanderligt maa danne Grundlaget for
alle Forandringer, tager Videnskaben sin Begyndelse. Opgaven er, at
undersøge Forholdet mellem Tingene og Lovene. Tingene sandses, Lovene
tænkes; men da Tingene lige saa lidt kunne existere uden Love, som
Lovene kunne tænkes uden Ting, saa maa Sandsning og Tænkning oprindelig
forudsætte hinanden i Fornuftens Eenhed.

At den religiøse og den videnskabelige Erkjendelse maae ifølge den
angivne Charakteristik være væsenlig forskjellige, er da vel indlysende.
Den religiøse Tro, der har sin Viden i Phantasi og Selvfølelse, kaster
et Blik paa Tingenes naturlige Verden, forvisser sig om dens
Forkrænkelighed og griber det Overnaturlige; tynget af Loven vender
den sig mod Miraklerne og seer i dem Vidnesbyrd om, at den frie,
almægtige Gud er Lovenes Herre. De religiøse Love ere Frihedslove,
Opdragelseslove, der kunne forandres efter Tid og Omstændigheder,
saasom Lovene om rene og urene Dyr; de videnskabelige Love,
mathematiske, mechaniske, physisk-chemiske o.\ s.\ v., høre derimod
saa væsenlig med i Tingenes Natur, at de umulig kunne forandres, uden
Naturen selv maa forandres. Alle Videnslove ere Nødvendighedslove;
kunde de brydes eller ophæves, saa \opage{7} maatte den hele Verden falde sammen;
Videnskaben maa desaarsag i Fornuftens Navn resolut angribe det
Mirakuløse.

Der er vel dem, der mene, at dette er en for haard Tale. Der er i den
virkelige Tilværelse ganske vist Knudepunkter, hvori det Overnaturlige
og det Naturlige, Miraklet og Loven synes at brydes med hinanden; men
det er kun tilsyneladende: det nye Reale, der indsættes ved Miraklet,
maa jo ved at indgaae i Sammenhæng med det Naturlige, rette sin
Udvikling efter Loven. Vistnok, svare vi, kan man i Almeensætningers
svævende Form forlige det indbyrdes Modstridende, men hvad hjælper det,
naar Modsigelsen viser sig i Kjendsgjerningen selv. Der berettes (2den
Moseb.\ III, 2--4) om Mose Kaldelse: „Jehovas Engel lod sig tilsyne
for ham i en Ildslue midt ud af en Tornebusk, og han saae, og see
Tornebusken brændte i Ilden, og Tornebusken blev ikke fortæret. Da
sagde Moses: jeg vil vige til Siden og see det store Syn, hvorfor
Tornebusken ikke opbrændes. Men Jehova saae, at han veg tilside for at
see, og Gud kaldte paa ham midt ud af Tornebusken og sagde: Moses!
Moses! og han sagde: see, her er jeg“. Hvorledes det vil være Troen
muligt at opfatte denne Beskrivelse, kunne vi endda nok forstaae, thi
Troen seer Alt i Phantasiens Lys, og Phantasien veed, hvorledes det
Overnaturlige forbinder sig med det Fornuftige, Lovenes Gyldighed med
Lovenes Afbrydelse; men i Viden og Videnskab lader det sig ikke
opfatte. Den menneskelige Viden er i Alt, baade Sandseligt og
Tænkeligt, saa strengt bundet til det Naturlige, det blot Fornuftige,
at den hverken kan høre Gud tale som et Menneske ei heller indsee,
hvorledes en Tornebusk kan vedblive at brænde uden at fortæres af
Ilden. Heraf er det da klart, at et Troesobject maa være af en ganske
anden Natur end et Vidensobject. Stille vi to Naturforskere, een paa
hver Side af Moses, de vilde da kunne see Tornebusken, saa sandt det er
en udvortes virkelig Gjenstand, men de vilde umulig enten kunne see
Engelen eller høre Jehovas Ord. Tornebusken er et Vidensobject, men
Engelen og Jehova ere Troesobjecter.
Hvad kan nu heraf udledes? At Videnserkjendelsen, Erkjendelsen af det
Naturlige, maa være \opage{8} fælles for alle naturlig Begavede, der gjøre
fornuftig Brug a % sic: the print has „Brug a“ (dropped f), p. 8
Sands og Forstand, medens Troeserkjendelsen,
Erkjendelsen af det Overnaturlige, saa sandt en saadan Erkjendelse
virkelig gives, kun er mulig for „de Udvalgte“, for dem, som med særlig
Sands for det Overnaturlige kunne see, hvad Andre ikke kunne see, og
høre, hvad Andre ikke kunne høre.

% --- p. 8 ---
Vi indsee, at der umulig kan gives en Erkjendelse, uden at
der maa være Nogen, som erkjender, og Noget, som bliver erkjendt,
skjøndt Forholdet er et ganske Andet i Religionen end i Viden\-%
skaben. At enhver Erkjendelse af det Objective maa forudsætte
en Objectivering, og enhver Objectivering en objectiverende Sub\-%
jectivitet, vil vist ingen benægte, som har tænkt alvorlig over
Sagen. Der kan være utallige Objecter, som endnu intet Menneske
har objectiveret, og der kan maaskee gives Objecter, til hvis Ob\-%
jectivering det mangler de nødvendige Betingelser; men Objecter
saa objective, at de hverken lade sig erfare eller forestille eller
tænke af nogen som helst Subjectivitet, kan der ikke gives. Men
staaer det urokkelig fast, at intet Object kan være uden en til\-%
svarende Objectivering, saa følger jo uimodsigelig heraf, at en
objectiverende Subjectivitet maa strække sig lige saa vidt som
Objecternes Verden. Omvendt, kan Subjectiviteten være saa uende\-%
lig klar, saa evig og dyb, man vil, den maa dog svinde som et
tomt Skin, naar den sidste Rest af det Objective forsvinder. Sub\-%
jectiviteten er lige saa lidt i Stand til at objectivere sig selv uden
at forudsætte et Andet, der ikke er Subjectivitet, som det saakaldte
rene Rum eller den rene Tid vil have Bestaaen, naar det virkelige,
foranderlige Indhold tænkes at være borte. Principielt ere Sub\-%
jectivt og Objectivt modsatte og dog sammenhørende Bestemmelser,
der staae og falde med hinanden. At tale om det for Mennesket
Objective er, saavel i Religionens som i Videnskabens Sprog, at
tale om det, som menneskelig lader sig objectivere. Objecternes
Forskjel er væsenlig betinget af en forskjellig Objectivering, og
Objectiveringens Forskjel er begrundet i den objectiverende Sub\-%
jectivitet.

% --- p. 9 ---
\opage{9}Aabenbaringen, Religionens Objective, kan kun objectiveres i
Phantasi. Hermed paastaaes ingenlunde, at det Aabenbarede kun
er en Phantasi. At fastholde en saadan Paastand vilde i sig selv
være lige saa urimeligt som at paastaae, at et Object som en
Steen, en Klippe kun var en Sandsning, fordi slige Gjenstande
nødvendig maae objectiveres ved Sandsning. Vi betragte ikke her
Phantasien som skabende Evne, men som Erkjendelsesmedium.
% saa vel: a comma-like speck stands before "vel" in the scan; not transcribed
Den er nødvendig saa vel i Videnskaben som i Religionen,
men dens Brug er lige saa forskjellig, som den videnskabelige
Objectivering er forskjellig fra den religiøse. I Videnskaben falder
Interessen væsenlig paa det Objective, i Religionen derimod først
og sidst paa det Subjective. En videnskabelig Objectivering f.\ Ex.
af Solsystemet forudsætter Phantasi, men Phantasien faaer des\-%
uagtet ikke Lov til at indføre et eneste Led, end ikke et eneste
Træk uden at dets Realitet er tilstrækkelig prøvet i Erfaring og
Tænkning. I hvor vel al Viden, som vi have seet, nødvendig
% sic: comma after "bestemt" as printed
maa forudsætte en for det Objective bestemt, Subjectivitet, saa er det
dog i Videnskaben en uafviselig Fordring, at Objecterne ikke
farves af subjectivt Lys, men uden al Indblanding af vilkaarlige
Elementer eller uvedkommende Tilsætninger opfattes som de maae
antages at være i sig selv efter almeengyldige Love. Hvorvidt
slige Objecter ere Mennesket behagelige eller ubehagelige, gavn\-%
lige eller skadelige, faaer ingen Indflydelse paa den videnskabelige
Objectivering; Videnskabens Objecter ere i denne Forstand uper\-%
sonlige. Anderledes med de religiøse Objecter. I Religionen
maa den objectiverende Subjectivitet tage Hensyn til sit eget
% sic: "Frygt." with a full stop (no comma tail at 600 ppi) as printed
Indre, sin Frygt. sit Haab, sin Trang og Lidenskab, sin Fred og
Frelse. Objectiveringen sætter da Sindet i Bevægelse, og Phan\-%
tasien meddeler Objecterne et personligt Anstrøg. I Davids
Psalmer er dette Træk paa mange Maader fremtrædende. I Psalm.
XVIII. 7--18 hedder det: „I min Angst paakaldte jeg Je\-%
hova \ldots\ Da blev Jorden rystet, saa den bævede og Bjergenes
Grundlag skjalv \ldots\ Han red paa en Cherub og fløi, han svævede
paa Veirets Vinger. Han slog et Telt af Mørke som sit Læ op
% --- p. 10 ---
\opage{10}omkring sig. Det var hans Hytte, mørke Vande i tætte Skyer“.
Her objectiveres i Phantasi. Jvnfr.\ Ps.\ LXV, Ps.\ CIV, Ps.\ CXIV o.\ s.\ v.

At Religionen, om end i anden og høiere Forstand, lige saa
vist behøver Phantasiens frie Brug til sin Objectivering af Virke\-%
ligheden som Poesien, kan blandt Andet sees af Skabelseshistorien.
Naar det i 1.\ Moseb.\ 1.\ Cap.\ hedder, at Guds Aand svævede over
Vandene, og at Gud sagde bliv Lys, saa bliver Aanden synlig i
Phantasiens Medium, og Gud selv viser sig som et Troens Object,
idet han frembringer Lyset. Ved Lyset og dets Modsætning
Mørket, maa vel menes det sandselige Lys og det sandselige
Mørke --- „det blev Aften og det blev Morgen“ --- men det Sand\-%
selige opløser sig i det Oversandselige, thi Gud er ingen sandselig
Gjenstand, og Ordet: bliv Lys! kan kun høres med Troens Øre.
Den videnskabelige Lyslære kan ikke opdage det Lys, der skabes
paa den første Dag og gaaer forud for de lysende Legemer, det
store og det lille Lys, der først blive til paa den fjerde Dag.
Heller ikke kan Videnskaben finde sig tilrette med, at der af Jor\-%
den fremspire Urter en heel Dag før Solens, Maanens og Stjer\-%
nernes Tilblivelse, saa lidt som den kan faae Øie paa den ud\-%
strakte Befæstning, der skiller mellem Vandene ovenfor og nedenfor.
Har Troen Ret, saa maa Videnskaben have Uret. En Bestræbelse
for at forlige Tro og Viden ved at lade de sex Dage betegne sex
Perioder, hver paa flere Millioner Aar, er lige ødelæggende for
Religion og Videnskab. Ved at gaae ind herpaa kommer Reli\-%
gionen jo bort fra Bibelens klare Ord om Aften og Morgen, og
finder ingen Forstaaelse af 2det Cap.\ 3.\ V., hvor Gud helliger den
syvende Dag fordi han hvilede paa den fra al sin Gjerning. Vi\-%
denskaben kan heller ikke være tjent med en saadan Overeens\-%
komst; den mangler Øie for det Overnaturlige, og vil umulig
kunne finde sig tilrette med, at Urter og Træer ere fremkomne
Millioner Aar, før Solen blev til.

At den religiøse Objectivering af Tingene og deres Sammen\-%
hæng er grundforskjellig fra den videnskabelige, maa Enhver, som
vil sætte sig ind i Forholdet, nødvendig indrømme. Betragte vi
Paradiset, Kvindens Skabelse og Syndefaldet, 1.\ Moseb.\ II.\ og III.,
% --- p. 11 ---
\opage{11}vil Forskjellen mellem Tro og Viden vise sig, om muligt,
endnu mere iøinefaldende. For Troens Objectivering har Eden,
den Have, som Gud plantede i Øst, og hvori han satte Menne\-%
sket, ikke blot lige saa megen Virkelighed som enhver anden Have,
men uendelig mere; thi her smelter det Naturlige harmonisk
sammen med det Overnaturlige. Floden, der deler sig, Guldet,
Bdellion og den Steen Schoham ere jo naturlige, Livets Træ
overnaturligt, og Kundskabens Træ paa een Gang baade naturligt
og overnaturligt. Eller skulde der maaske slet intet Naturligt
være i Kundskabens Træ, hvorledes er det da muligt, at Eva, som
dog ikke kan være aldeles overnaturlig, kan forlystes ved at see
det og spise af dets Frugt. Men hvis Træet ikke tillige var over\-%
naturligt, hvorledes skulde det da foruden Kundskabens Frugt
kunne bære den talende Slange. At han, som, Joh.\ VIII, 44, kaldes
Manddraberen fra Begyndelsen af, maa have været i Slangen, er
ogsaa mere overnaturligt end naturligt. Vel har man ikke sjælden
erklæret, at den hele Fremstilling af Paradiset og Syndefaldet ikke
nødvendig maatte forstaaes bogstavelig, men meget vel kunde op\-%
fattes enten allegorisk eller mythisk, uden at den religiøse Tro
derved led noget Afbræk. At dette imidlertid er en for Aaben\-%
baringen farlig Indrømmelse, vil let falde i Øinene. Det kan
ikke nok som indskærpes, at al Aabenbaringstro i sit Princip er
og maa være Mirakeltro. Skulde det Mirakuløse i Paradiis vir\-%
kelig berettige os til at opfatte den hele Fremstilling som en
Mythe, da maatte det Samme jo gjælde om alt det Miraku\-%
løse saa vel i det Nye som i det Gamle Testamente. Christi
Fødsel af Jomfruen, hans Mirakler, hans Opstandelse og Himmel\-%
fart maa da opløse sig i lutter Myther. Men naar Miraklet, i
hvilken Skikkelse det end fremtræder, er Mythe og Aabenbaringen
Mythologie, saa er Religionen Videnskabens Dom underlagt, og
% sic: comma after "opløst" before capital "Staaer" (comma tail clear at 600 ppi)
dens positive Form dermed væsenlig opløst, Staaer man altsaa
stivt paa den Grundsætning, at vor Sandhedserkjendelse kun kan
have eet Princip, efterdi den evige, absolute Sandhed kun er een
og maa i Alt være sig selv lig: saa er Striden mellem Religion
% --- p. 12 ---
\opage{12}og Videnskab, mellem Tro og Viden, en Strid paa Liv og Død,
saa indtrængende og afgjørende som muligt.

Er det virkelig sandt, hvad Aabenbaringstroen lærer, at
Menneskehedens Historie gaaer ud fra et Paradiis, da kan det
umulig være sandt, hvad end Naturhistorien paastaaer, at Menne\-%
sket er blevet til ved at udvikle sig af en eller anden Dyreform;
er det virkelig sandt, at Christus er opstanden af Graven og le\-%
gemlig opfaren --- ikke til en anden Klode i Rummet, men til
Himmelen, saa maa der et Steds være en virkelig og dog over\-%
virkelig, naturlig og dog overnaturlig Himmel, hvori han boer, en
Himmel, som Astronomerne ikke kjende. Skjøndt al vor Erkjen\-%
delse er Stykværk, kan der dog ikke være Tale om at udstykke
Sandheden selv, saa at nogle sjældne, særdeles udmærkede Stykker,
som f.\ Ex.\ hvor Tyngdeloven ophæves og der vandres paa Havet,
eller hvor de chemiske Love sprænges, og Vand ved et Ord for\-%
vandles til Viin, høre under Religionen, medens de mere sædvan\-%
lige og naturlige Stykker hjemfalde til Videnskaben. Adskillelsen
gjælder ikke her de enkelte, snart overraskende, snart mere sæd\-%
vanlige Tilsyneladelser af Sandheden; nei, her handles om et
Princip, Principet for den menneskelige Erkjendelse af Sandheden.
Er her kun eet Princip, da maa det nødvendig sættes enten i
Troen eller i Viden. Sættes det i Troen, saa maa Videnskaben
ubetinget bøie sig for Religionen; Videnskaben kjender ikke Gud
og kan, hvor langt den end seer tilbage, dog ikke naae til Ska\-%
belsens Begyndelse. Sætte vi Principet i Viden, saa maa Reli\-%
gionen underordne sig under Videnskaben. Man kan da sige til
den, der ikke formaaer at tilegne sig den videnskabelige Forsk\-%
nings store Resultater: Du kan i din Enfoldighed jo gjerne troe
paa Paradiis og Syndefald, paa Christi overnaturlige Undfangelse
og Himmelfart, det i disse Billeder og Forestillinger for det me\-%
nige Folk Vækkende og Trøstende; men bild dig aldrig ind, at
det virkelig er Sandhed. Sandheden er kun i Videnskaben, og
den, der ikke kjender Videnskaben, faaer aldrig at vide, hvad
Sandhed væsenlig er.

At der kun er eet Erkjendelsesprincip vil i praktisk Forstand
% --- p. 13 ---
\opage{13}have til Følge, at der kun kan gives een allerhøieste Sandheds\-%
autoritet, som da maa søges enten i Religionen eller i Viden\-%
skaben. Hvad heraf følger, har Historien viist. I Katholicismen
hele Middelalderen igjennem, ja selv i den gamle Protestantisme
til henimod Slutningen af det syttende Aarhundrede, var Principet
religiøst og den overgribende Autoritet paa Kirkens Side; vel var
en videnskabelig Formudvikling ogsaa i Bevægelse, men hvad
Sandhedsindholdet angaaer, maatte den ved Syndefaldet formørkede
Fornuft laane sit Lys fra Aabenbaringen, ære denne som en Tje\-%
nestepige (ancilla) sin Herskerinde og finde sig i sit Fangenskab
under „Troens Lydighed“. Skal Troesprincipet gjælde som det
eneste sande Erkjendelsesprincip, da er det fuldkommen conseqvent,
at den Videnskab, hvori Fornuftens Form bliver tvungen til at
udvikle Aabenbaringens Indhold, nemlig Theologien, maa ansees
for den høieste og egenlig sande Videnskab. Hvad andre viden\-%
% altsaa: a low speck after "altsaa" read as dirt, not a full stop
skabelige Forsøg angaaer, er Theologien altsaa i sin fulde Ret,
naar den i stridige Tilfælde indkalder dem for sin Domstol og
dømmer dem efter sine egne Love.

I det attende Aarhundrede begynder Forholdet at vende sig;
ikke blot i historisk-philologisk Kritik, men ogsaa i Philosophi og
Naturvidenskab, som trods Strid og Uenighed indbyrdes, vare enige
om at anerkjende Fornuftens Ret, emanciperede Videnskaben sig
fra Religionens Autoritet, erklærede reent ud, at det sande Er\-%
kjendelsesprincip ikke var religiøst, men videnskabeligt. Først
søgte man at udrense ethvert Spor af det Uvidenskabelige og for\-%
vandle Religionen selv til den høieste Videnskab; da man opda\-%
gede Feiltagelsen, lagde Videnskaben, i Kraft af det enegyldige
rationelle Princip, directe an paa en radical Opløsning af al positiv
Religion. De, som endnu i Aabenbaringens Navn ville gjøre Troen
paa det Overnaturlige til det eneste sande, altomfattende Erkjen\-%
delsesprincip, have mørke Udsigter; thi Videnskaben med For\-%
nuften og den naturlige Verden til Grundlag er en Magt, ja en
Stormagt, der ikke anden Gang lader sig kneble og lægge i Læn\-%
ker. Det maatte da skee ved et Mirakel saa stort, at hele den
nyere Kultur styrtede sammen i Chaos.

% --- p. 14 ---
\opage{14}Den protestantiske Kirke har, hvad Miraklerne angaaer,
neppe været heldig med sin Bibelfortolkning. I Marcus' Evan\-%
gelium XVI. 15--18, siger den himmelfarende Frelser: „Gaaer
ud i al Verden og prædiker Evangeliet for al Skabningen. Hvo
som troer og bliver døbt, han skal vorde frelst, men hvo som
bliver vantro, han skal vorde fordømt. Men disse Jertegn skulle
ledsage de Troende: i mit Navn skulle de uddrive Djævle, de
skulle tale med ny Tunger; De skulle tage paa Slanger, og der\-%
som de drikke nogen Forgift, skal det ikke skade dem; de skulle
lægge deres Hænder paa de Syge og de skulle helbredes“. Ingen
vil kunne nægte, at den almindelige Kirke fra sit Synspunkt har
Ret, naar den i ubetinget Ærefrygt for Christi Ord lægger afgjø\-%
rende Vægt paa Troen og Daaben som ufravigelige Betingelser
for Frelse og Salighed; men det maa i høieste Grad forundre os,
at de protestantiske Kirker, selv den bibelfaste, orthodoxe Luther\-%
dom, aldeles intet Hensyn tager til de udvortes Kjendetegn, hvor\-%
med Christus selv i Afskedens Øieblik høitidelig betegner de Tro\-%
ende. Det er en Polemik imod Katholicismen, der har forledet
den protestantiske Kirke til at miskjende den uopløselige Forbin\-%
delse mellem den bibelske Christendom og dens Mirakler. Man
forarges over de katholske Helgenlegender, de mange smaalige, smag\-%
løse Mirakelfortællinger, og griber saa til den, som det synes, meget
fornuftige Udvei, at lade Mirakeltiden ende med Apostlenes Død
og den rene Christendom fortsætte sig, ved Ordets og Sakramen\-%
tets hemmelighedsfulde Virkninger alene, uden synlige Mirakler.
I Holbergs Kirke-Historie, 1ste Deel\footnote{Udgivet ved Liebenberg. 1867.} % *)
S. 31 hedder det om Chri\-%
stendommen: „Ingen Lærdom har været bekræfted med saa mange
og store Miracler, skiønt ingen Lærdom behøvede mindre Jærtegn,
og det i henseende til dens Ypperlighed; thi alt hvad Jesus lærede
paa Jorden, sigter allene til det Menneskelige Kiøns Lyksalighed,
og seer man, at dets Morale har haft saadan Virkning, at Barba\-%
riske Lande endogsaa paa de Tider, da Lærdommen ved Menne\-%
skelige Paafund var bleven fordunkled, deraf have faaet en anden
% --- p. 15 ---
\opage{15}Skikkelse. Dog haver Gud ligesom til Overflod authoriseret den
ikke alene ved Propheternes Spaadomme, men ved utallige Miracler,
paa det at ingens Vantroe kunde undskyldes“. --- Men hvo seer nu
ikke, at en saadan Skilsmisse, som Holberg ingenlunde har op\-%
fundet, men udtrykkelig forefundet i den lutherske Kirke, Skilsmissen
mellem en Christendom med Mirakler og en Christendom uden
Mirakler, er noget i det Nye Testament saa aldeles Ubekjendt, at
hverken Christus eller Apostlene synes at have endog den fjerneste
Anelse derom. Ved at gaae ind paa denne Skilsmisse har den
lutherske Kirke, trods alle Protester, uden at være sig Modsigelsen
bevidst, forflyttet Erkjendelsesprincipet fra Religionen over i Vi\-%
denskaben og derved gravet sin egen Grav. Det er da uimodsi\-%
gelig vist, at Aabenbaring og Fornuft, Tro og Viden, maae føre en
uophørlig, uforligelig Strid med hinanden, indtil enten Religionen
har underlagt sig al Videnskab eller omvendt Videnskaben gjort
sig til Herre over al Religion. Kan man umulig gaae ind
paa den første Antagelse uden at krænke Videnskaben, og umulig
give sig hen til den sidste Opfattelse uden at krænke Religionen
og angribe Folkelivet i dets inderste Rod, saa bliver det nødven\-%
digt at undersøge, om ikke den ene, evige, i sig selv absolute
Sandhed skulde ved sin Afspeiling i Menneskets Bevidsthed
være nødt til at dele sig selv, straale fra to forskjellige Centra
og derved afsætte to relative, dog for al menneskelig Erkjendelse
absolut ueensartede Principer, det ene for Troen, det andet
for Viden.

At de to Sandhedskredse, hver med sit Centrum, ikke kunne
staae i Strid med hinanden, følger deraf, at begge have deres
Udspring fra den ene og samme evige Sandheds-Eenhed; men at
de heller ikke kunne smelte sammen, følger deraf, at deres Er\-%
kjendelse, paa Grund af Adskillelsen, maa have forskjelligt Ud\-%
gangspunkt, forskjelligt Maal og forskjellig Methode. Troesobjectet
kan umulig blive Vidensobject, ei heller omvendt. Spørgsmaal
som de om Skabelsen, Paradiset og Syndefaldet, Forløsningen og
Gjenopreisningen kunne slet ikke besvares videnskabelig, og maae
desaarsag resolut afvises i Videnskaben; hvorimod alle Spørgsmaal
% --- p. 16 ---
\opage{16}om Atomer, Kræfter, Love, om Tingenes naturlige Dannelse og
Opløsning høre Videnskaben til og aldeles ikke vedkomme Reli\-%
gionen. Naar en Religionslærer rykker ud med en formeentlig af\-%
gjørende Polemik mod Videnskaben og vil, at dens Kritik af den
bibelske Skabelseslære, Miraklerne o.\ s.\ v. skal bevise, at Viden\-%
skaben er falsk, da beviser han kun, at den videnskabelige Er\-%
kjendelse hverken er ei heller vil eller kan være religiøs, just fordi
den er og maa være videnskabelig. Naar omvendt Nogen i Vi\-%
denskabens Navn beviser, at saadanne Forestillinger, som de om
Paradiis og Syndefald, Opstandelse fra de Døde o.\ s.\ v., fuld\-%
stændig maae udryddes, fordi de alle ere behæftede med utaalelige
Modsigelser: da have slige Beviser udelukkende Betydning fra et
videnskabeligt, ikke fra et religiøst Synspunkt. Religionen kan
med sit Princip, naar den taler reent ud, aabent og frit erklære,
at den er og maa være uvidenskabelig. Hverken Abraham eller
Moses eller Christus have været Videnskabsmænd eller talt ud af
dybere videnskabelige Forudsætninger; deres Grunde ere personlige,
og Erkjendelsens Objecter, der sees i Personlighedens Lys, faae
derved alle uden Undtagelse et personligt Anstrøg. Den Religiøse
kan paa sin Viis tage alle menneskelige Kræfter i Beslag uden
at gjøre Brug af Videnskaben; der gaaer igjennem det Gamle og
Nye Testamente en rig og dyb og mægtig Strøm af Tanker, men
de have væsenlig Intet med Videnskaben at gjøre. Videnskaben
kan efter sit Princip frit ud erklære, at den i Grunden hverken
er christelig eller jødisk eller hedensk, men almeenfornuftig og
maa i det Almeenfornuftige ene og alene søge sin oprindelige
Selvstændighed.

Først naar denne Grundadskillelse, denne Timelighedsdualisme,
hvori det Evige indgaaer, tilfulde erkjendes og fastholdes, kan der
være Tale om sand Vexelvirkning mellem det Religiøse og det
Videnskabelige i Menneskelivet. Den evige, absolute Sandhed,
der i sit Væsen er monistisk, gaaer ind i Timeligheden og bliver
dualistisk; Mennesket, i hvem denne Dualisme kommer til klar
% asketisk-uvidenskabelig: line-end hyphen kept as a compound hyphen (cf. nihilistisk-videnskabelig, p. 17)
Bevidsthed, kan ikke blive Dualismen kvit enten ved asketisk-%
uvidenskabelig Fornægtelse af det Timeliges Realitet eller ved
% --- p. 17 ---
\opage{17}nihilistisk-videnskabelig Fornægtelse af det Eviges Realitet, men maa
i dualistisk Erkjendelse af begge Realiteter arbeide sig frem i Til\-%
nærmelse til stedse rigere Forstaaelse af den absolut sande, con\-%
crete Monisme.

Ved i Korthed at belyse de historiske Misforhold mellem
Religion og Videnskab skulle vi i det Følgende forsøge at arbeide
os hen til et Standpunkt, hvorfra det kan sees, at Religionen lige
saa lidt kan være Videnskabens som Videnskaben Religionens For\-%
mynder, men at begge, netop i Kraft af de to ueensartede Prin\-%
ciper, kunne arbeide sammen i god Forstaaelse.

% print: a short centred rule separates Section I from the head; head numeral in heavier type
\begin{center}II.\end{center}

Er den aabenbarede Religion Sandheden selv, og Sandheds\-%
erkjendelsen kun een, saa har Fornuften intet særligt Erkjendelses\-%
princip. Aabenbaringen maa da være uendelig ophøiet over Men\-%
neskets indskrænkede, ved Synden formørkede Fornuft; ei engang
Forklaringen af den naturlige Tingenes Orden tør overlades en for\-%
nuftig Brug af Forstanden og de fem Sandser alene, thi Sandserne
gaae jo kun paa det Lavere og den til Fornuften overladte Tænk\-%
ning er udsat for store Vildfarelser. Luther, der ellers sætter
Fornuften høit i alle verdslige, borgerlige Sager, godtgjør dens
fuldstændige Afmagt og Utilstrækkelighed, naar Talen er om en
Besvarelse af afgjørende religiøse Spørgsmaal; navnlig har han i
Nadverstriden saaledes affeiet den Carlstadske Fornuft, at han
derved følelig har ramt den almeenmenneskelige. Man skulde da
mene, at den lutherske Kirke, om den end ved passende Brug af
en underordnet Fornuft kunde faae en vel ordnet og sammenhæn\-%
gende Religionslære, umulig kunde opnaae at faae en Dogmatik
saa videnskabelig alsidig begrundet som den scholastisk katholske.
Og dog er det skeet. Lutheranernes Troesbekjendelse har i den
Grad udviklet sig til Videnskab, at Dogmatiken er bleven afgjø\-%
% --- p. 18 ---
\opage{18}rende Norm for sand Rettroenhed. Melanchthon, der med sine Loci
Theologici begyndte religiøst og endte videnskabeligt, gik i Spidsen;
men da han ikke ret kunde faae Bugt med Fornuften og navnlig
ikke faae Øiet op for, at det Høieste netop maatte være det mod
Fornuften Stridende, blev hans Retning alvorlig bekæmpet. Først
i Formula Concordiæ fastsættes det udtrykkelig, hvorledes Fornuften
ubetinget er underkastet Orthodoxien, skal ydmyge sig under det
Overfornuftige, ja bøie sig dybt for det Modfornuftige, og nu blev
Religionslæren orthodox systematisk. Religionen blev atter Viden\-%
skab, idet den blev Theologi. Som en Ørn paa mægtige Vinger
hævede sig den gjenfødte Theologi i Aabenbaringens Sollys til en
saadan Høide, at den frit kunde see ned paa og bedømme al den
menneskelige Viden, der maatte holde sig ved Jorden. Hvorledes
det egenlig forholder sig med denne orthodoxe Videnskab, skulle
vi i faa Træk søge at belyse.

Som Erkjendelsens altomfattende Princip har Troen den hellige
Skrift til Grundlag; det gjælder da først om at sikkre sig Skriftens
Ufeilbarhed ved at sikkre sig dens verbale Inspiration. Det maatte
bevises, og blev da, om ikke egenlig beviist, saa dog i Troen paastaaet,
at Propheter og Apostle ikke af naturlig Drift eller paa egen Haand
have skrevet et eneste Ord, men at det udtrykkelig er skeet ifølge
den Hellig Aands umiddelbare Indskydelse. Hvert Bibelens Ord
er da i egenlig Forstand Guds Ord; de til Nedskrivningen udvalgte
Redskaber ere „Guds Skrivere, Guds Penne“. At Skriftens Auto\-%
ritet ikke, hvad Katholikerne paastaae, afhænger af Kirkens Auto\-%
ritet, men alene af Gud, kan sees af mangfoldige Skriftsteder, Joh.
V, 34; VII, 17; 1 Thess.\ II, 13 o fl.\footnote{See f.\ Ex.\ Brochmand.
Universæ Theologiæ Systema. Articul II. Cap.\ II p.\ 14--15.}. % *)
Hensigten med, at Gud
selv har villet meddele os sin Villie skriftlig, er alene den, at vi
skulle lære at dyrke ham paa rette Maade og føre vort Liv saa\-%
ledes, at vi sikkert kunne naae Saligheden\footnote{Anfrt.\ Skr.\
Cap.\ I. S.\ 2.}. % **)
Skriftens Klarhed,
Fuldstændighed og Nøiagtighed er saa stor, at selv Vocalerne i
det Gamle Testamente maae antages at være inspirerede. At der
% --- p. 19 ---
\opage{19}umulig kan paavises Modsigelser mellem forskjellige Skriftsteder,
hvor der udtrykkelig tales om det Samme, er en Selvfølge. Det
kan unægtelig gaae saa vidt, at Jehova paa et Sted udtrykkelig
siges at gjøre, hvad Satan gjør paa et andet, som naar det i 2den
Samuel XXIV, 1 berettes, at Jehova selv fristede David, medens
det 1ste Krøn.\ XXI, 1 hedder, at det var Satan, der fristede
ham; men dette kan meget godt forklares, naar man kun vil tage
Fornuften fangen. At Jehova frister, vil jo sige, han kan ikke
friste, --- „Ingen sige, naar han fristes: jeg fristes af Gud; thi
Gud fristes ikke af det Onde, og han frister selv Ingen“ (Jakobs
Br.\ I, 13). Da nu Jehova i sin Vrede vil friste David, er det jo
ikke uforstaaeligt, at han lader Satan gjøre, hvad han ikke selv
kan; kun vogte man sig for den Fornuftens Spidsfindighed, at lade
Jehova med Hensyn til Fristelsen selv være enig med Satan.
Heller ikke, hvad det reent Historiske angaaer, kunne de hellige
Forfattere, hvad dog Socinianerne paastaae, antages her og der at
have feilet. Rigtignok er der selv i det Nye Testamente adskillige
vanskelige Steder. Det er jo virkelig et stort Spørgsmaal, angaa\-%
ende Luc.\ II, 1--2, om Qvirinus netop var Landsherre i Syrien,
da Keiser August lod Alverden beskrive til Skat; det er, om mu\-%
ligt et endnu vanskeligere Spørgsmaal, hvorledes Propheten Jere\-%
mias, Matth.\ XXVII, 9, kan ombyttes med Zacharias. Kan man
ikke tage Fornuften til Fange, maa man indrømme en Hukommel\-%
sesfeil; men naar Troen er blevet Fornuftens Herre, indseer man
strax, at Jeremias jo meget vel kan have sagt, hvad ikke han selv,
men Zacharias har skrevet\footnote{Brochmand. Systema, 1ste Deel,
S.\ 81: Nobis tutissima responsio est hæc: verba illa qvæ a Matth.\
adscribuntur Jeremiæ, a Jeremiæ qvidem dicta esse, (non enim dicit
% sic: "al Jeremiæ" as printed
Evangelista, hoc al Jeremiæ scriptum, sed
% Greek fount in the print (the only Greek word); needs Greek font support to compile
ῥηθὲν ceu dictum est): verum a
% sic: "Zaeharia" as printed
Zaeharia in litteras sacras relata. Nam aliæ qvæ afferri solent
% sic: "græc-cum" broken over the line as printed
responsiones, aut Textum græccum corruptelæ arguunt, aut Matthæo
dicam lapsus memoriæ scribunt, aut conjecturis plane incertis
nituntur.}. % *)

Da Bibelen imidlertid underligger Fortolkning, og enhver
Fortolker drives af sin Aand, saa opstaaer det alvorlige Spørgs\-%
% (word continues on p. 20: "maal")
% --- p. 20 ---
\opage{20}maal: paa hvilken Fortolkning skal man da kjende den sande
Aand; Luthers Erklæring: „Der heilige Geist ist der allerein\-%
fältigste Schreiber und Redner, der im Himmel und Erden ist,
darum auch seine Worte nicht mehr denn einen einfältigen Sinn
haben können, welchen wir den schriftlichen oder buchstablichen
Zungen-sinn nennen“\footnote{Walch XVIII 1602.} % *)
giver os en foreløbig Veiledning. Kun i
Bibelens Aand kan Bibelen fortolkes; Bibelens Aand er den Hellig
Aand, følgelig maa Fortolkeren have den samme Hellig Aand,
hvis hans Fortolkning skal være sand. Da Jesus blev døbt af
Johannes, hedder det: „da aabnedes Himlene over ham“, idet han
steg op af Vandet, „og han saae Guds Aand fare ned som en Due
og komme over ham“. Og see fra Himlene kom der en Røst,
% sic: no opening mark for the quotation closed after "Velbehag"; the page has one more closing than opening mark here
som sagde: denne er min Søn den elskelige, i hvem jeg haver
Velbehag“. Matth.\ III, 16--17; hos Luc.\ III, 22 siges det udtryk\-%
kelig, at den Hellig Aand foer ned i legemlig Skikkelse som en
% Greek fount in the print. sic: "πνῦμα" (for πνεῦμα), unaccented "το", and "ἔιδει" (for εἴδει), all as printed
Due (καταβῆναι τὸ πνῦμα το ἄγιον σωματικῷ ἔιδει). Skal man
nu tage dette ganske eenfoldig efter Bogstaven, da er en klar og
bestemt Forestilling om, hvorledes Aanden har taget sig ud i en
virkelig Due-Skikkelse vanskelig, thi den, der seer eenfoldig paa
Duen, kan ikke see Aanden, og den, der fæster sit Blik paa Aan\-%
den, kan ikke see Duen; lige saa vanskelig bliver det at forstaae,
om det alene var Jesus, der saae Duen og hørte Røsten, eller
tillige Johannes. Naar man læser Johannes' Evangelium I, 15--32,
III, 27--36, saa skulde man jo troe, at Døberen selv baade havde
seet Duen og hørt Røsten; men sammenligne vi hermed Matth.\ %
XI, 2 og 3, og Luc.\ VII, 18--35, saa fristes man til at betvivle,
at Døberen virkelig har været enig med sig selv om, hvad han
egenlig har seet og hørt.

Forholdet mellem Aand og Bogstav er ingenlunde saa fat\-%
teligt, som det ved første Øiekast kunde synes. Begynde vi med
den bogstavelige Mening, saa faaer Fornuft og Forstand saameget
at gjøre med Ordbetydning og Sammenhæng, med Grammatik og
Hermeneutik, at Exegesen, videnskabelig tagen, let forglemmer den
% --- p. 21 ---
\opage{21}Hellig Aand; og dog er det Aanden, der giver Forstaaelsen sit
rette Præg. At der nu er Aand i den katholske Kirke lige saa
vel som i de protestantiske, vil dog vel Ingen nægte. Disse Kirker
ligge netop i Strid med hinanden, fordi enhver af dem paastaaer
at have Aabenbaringens sande Aand. Katholiker, Zwinglianere,
Calvinister vindicere sig hver især at have den Hellig Aand frem\-%
for Lutheranerne; hvorledes skal denne Strid nu afgjøres ved at
paakalde Aanden? Vi høre Luther ivre mod de falske Propheter:
„Schrift, Schrift, Schrift aber ist nicht eitel Geist, davon sie geifern,
der Geist müsste es allein thun; die Schrift sei ein todter Buch\-%
stabe, und könne nicht das Leben geben. Rühme aber nicht viel
vom Geiste, wenn du nicht das aeusserliche offenbare Wort hast,
denn es wird gewisslich nicht ein guter Geist sein, sondern der
leidige Teufel aus der Hölle“. Vi kunne heri give Luther Ret,
og komme dog ikke et Haarsbred videre; hans Ord til Zwingli\-%
anerne paa Colloqviet i Marburg 1529: „Ihr habt einen andern
Geist denn wir“, viser tydelig nok, hvad den paaberaabte Eenhed af
Legeme og Brød, af Blod og Viin, vil sige. Skal det Afgjørende ligge
i Troen, saa maa Enhver jo blive ved sin Tro; Zwinglianeren har
da selv samme Ret som Lutheraneren. Om end Indstiftelsens
Ord: „dette er mit Legeme, dette er mit Blod“, tages og bruges
ganske efter Bogstaven, saa er der dog i Bogstaven selv ingen
Forstaaelse; Forstaaelsen, der kommer med Udtydningen, maa nød\-%
vendig blive enten fornuftig: „dette betyder“, eller modfornuftig:
„dette \emph{er}, hvad det dog endnu ikke, hverken naturlig eller over\-%
naturlig, kan være“. Da Christus uddeelte Brød og Viin, havde
han jo endnu sit eget naturlige Legeme og Blod; skal Brødet og
Vinen \emph{ikke blot betyde}, men paa \emph{overnaturlig Maade
være} det naturlige Legeme, der skal hengives, det naturlige Blod,
der skal udgydes paa Korset, saa maae vi nødvendig enten faae
en aldeles uforstaaelig Tvetydighed, eller et Mysterium saa uigjen\-%
nemtrængeligt, at Troen selv maa klamre sig til det Absurde ---
credo qvia absurdum est. Men det Absurde giver ikke blot ingen
videnskabelig, men heller ingen religiøs Forstaaelse; det standser
enhver mulig Opfattelse ved det absolut Meningsløse. Til en
% --- p. 22 ---
\opage{22}saadan Conseqvens har den lutherske Kirke ikke kunnet hengive
sig; den er bleven staaende ved Troen, d.\ e.\ ved \emph{sin} Tro paa
Mysteriet.

Maa nu hvert særeget Kirkesamfund være sig bevidst at
have en særegen Tro, saa maa det jo ogsaa have særegne Formler
for sin eiendommelige Troesbekjendelse, særegne Symboler. Naar
da Kirketroen ene og alene skal bygges paa den hellige Skrift,
saa maa Symbolernes Gyldighed nødvendig beroe paa, at de nøi\-%
agtig stemme med Skriftlæren. Den lutherske Kirke paastaaer
ikke, at dens Symboler ere ufeilbare, den holder kun paa dem
\emph{fordi} (qvia) og \emph{for saa vidt} (qvatenus) deres Indhold er i væ\-%
senlig Overeensstemmelse med Skriftens Indhold. Men her kommer
Fortolkningsspørgsmaalet igjen, thi Bevægelsen gaaer i en Cirkel.
En Fortolkning uden særegen Troesbekjendelse kan være fornuftig,
skarpsindig, dybsindig og dog ikke sand, fordi Sandhedens Aand
mangler; men Sandhedens Aand er Troens Aand, og den sande
Tro lever kun i den sande Kirke. For Lutheraneren maa den
lutherske Kirkes Aand være Eet med Aabenbaringens Hellig Aand;
Lutheraneren maa altsaa forstaae Symbolerne i Bibelens Aand,
men deraf følger omvendt, at han maa fortolke Bibelen i Symbo\-%
lernes Aand; afviger han fra denne Regel, d.\ v.\ s.\ bryder han
Cirklen, bliver han en Kjætter.

Kaste vi nu et Blik paa Dogmerne, da er det ikke det for
de enkelte Kirker Særegne, men det for dem alle fælles, det Al\-%
mindelige, hvorved der i denne Sammenhæng skal dvæles. Synde\-%
faldet bevirkes ved en Fristelse af Djævlen, Opreisningen skyldes
udelukkende Christi Fortjeneste. For at paavise Dogmatikkens
videnskabelige Charakteer i det Hele, behøve vi altsaa blot at
kaste et Blik paa Forholdet mellem Christus og Djævelen i den
lutherske Dogmatik. Hvorvidt Christus vilde være blevet incar\-%
% sic: "atiamsi" (for etiamsi) as printed
neret, hvis Mennesket ikke havde syndet (atiamsi homo non pec\-%
casset), et Spørgsmaal, som Thomas Aqvinos udførlig behandler\footnote{Gent.\ IV c.\ 53, 55. Jvnfr.\ Dr.\ Karl Werner. Der heilige Thomas von
Aqvino 1859. 2ter Band. S.\ 620.}, % *)
% --- p. 23 ---
\opage{23}har den lutherske Kirke ikke indladt sig paa at afgjøre\footnote{Brochmd.\ S.\ 725--26.}; % *)
men
derimod i stor Udførlighed undersøgt Forholdet mellem Christi
guddommelige og menneskelige Natur, og derved gjort det ret
tydeligt, hvorledes Han, som blev aabenbaret, for at „\emph{forstyrre
Djævelens Gjerninger}“, virkelig havde Magt til at fuldende
sin Mission. Skjøndt Djævelens Magt virkelig blev brudt, da
Christus aabenbarede sig, saa vedbliver Djævelens Rige dog til det
Sidste at modstaae Christi Rige. Luthers Tro paa en virkelig,
personlig Djævel var lige saa urokkelig som hans Tro paa en vir\-%
kelig, personlig Christus. Luther er enig med Tertullian i, at ikke
blot onde Tanker og Begjæringer, men ogsaa Misvæxt, Pest, Ilde\-%
brand og andre udvortes Ulykker have Djævelen til Ophavsmand.
I Cathechismus major og paa mange andre Steder søger Luther
paa den meest indtrængende Maade at indskærpe, hvor vedholdende,
alsidig og utrættelig Djævlenes og Dæmonernes Virksomhed er for
at vi kunne være paa vor Post. „Die Heiden“, siger han, „wissen
nicht, woher das Unglück so plötzlich kommt; aber wir wissen es,
das es eitel Teufels Arbeit ist, der hat solche Helleparten, Blei\-%
kugeln und Büchsen, solche Spiesse und Schwerter, damit er unter
uns schiesst, wirft und sticht, wenn Gott es ihm erlaubt \ldots“; og
paa et andet Sted: „Ein Christ soll wissen, dass er mitten unter
den Teufeln sitze, und dass ihm der Teufel näher sei den sein
Rock oder Hemde, ja näher denn seine eigene Haut, dass er rings
um uns her sei, und wir also stets mit ihm zu Haar liegen und
uns mit ihm schlagen müssen“\footnote{Jvnfr.\ Walch.\ X, S.\ 1234, XIII, S.\ 2550.}. % **)

Det er dog ikke Luthers personlige Autoritet, men Bibelens, navn\-%
lig det Nye Testaments Fundamentallære, der gjør det nødvendigt
for den lutherske Dogmatik at bygge, om vi saa maa sige, paa
Djævletroen og følge Arvefjendens Kamp mod Christi Rige indtil
Dommens Dag, da han selv med alle sine Tilhængere for evigt
skal bindes i Helvede. De lutherske Dogmatikere ere desaarsag
lige saa omhyggelige med at undersøge, hvorvidt Djævelen egenlig
% --- p. 24 ---
% footnote: "ll." as printed (lowercase-looking); rejected alternative "II."
\opage{24}kan gaae, og hvorpaa hans Magt beroer, som med at udvikle Troen
paa Christi Fortjeneste\footnote{See f.\ Ex.\ Brochmand ll.\ Tom.\ I p.\ 250. Qvæst.\ XIII. An Diabolus
cogitationes cordis novisse posse. Pg.\ 251. Qvæst.\ XIV. An Diabo\-%
lus vera miracula edere qvæat. Pg.\ 253. Qvæst.\ XV. An Diabolus
animas demortuorum in hanc vitam revocare, et substantiales rerum
formas immutare possit? etc.\ etc.}. % *)
Vi maae i dette Stykke give Strauss
Ret, naar han gjør gjældende, „at Ideen om Messias og hans Rige
lige saa lidt er mulig uden et Dæmonrige, med et personligt Over\-%
hoved til Modstykke, som en Magnets Nordpol uden Sydpol“, vi
maae indrømme ham, at „den fromme Indskrænkethed, som
frygtede for at tabe Christus, naar de tabte Djævelen, har seet
langt rigtigere end Schleiermacher med sit Postulat, at Troen paa
Djævelen paa ingen Maade turde opstilles som Betingelse for Troen
paa Christus“\footnote{Die christliche Glaubenslehre von Dr.\ D.\ F.\ Strauss. Zweiter Band.
1841. S.\ 15.}. % **)

Den orthodox lutherske Dogmatik, hvis Blomstringstid falder
i det syttende Aarhundrede, maa vel betragtes som et i sit Slags
alvorligt Forsøg paa at fremstille Bibelens uforvanskede Indhold
i streng videnskabelig Form, et i Sandhed storartet Forsøg; men
det var den bundne Fornuft, der, som Troens Slave, maatte udføre
det hele Arbeide; den blev sig da efterhaanden sine Kræfter be\-%
vidst, drog lidt efter lidt Erkjendelsesprincipet over til sig, og for\-%
beredte derved sin fuldstændige Emancipation. Djævelen, der med
sine Dæmoner især havde mishandlet den, skulde nu selv afgive
den Fodskammel, den traadte paa, da den begyndte at hæve sig
til sin hele Høide. En reformeert Theolog i Holland, Balthasar
Bekker, ivrig Tilhænger af den Cartesianske Philosophi, udgav et
meget opsigtsvækkende Skrift: „Den forhexede Verden“,\footnote{De betoverede weereld, Amsterd.\ 1691--94. Paa Tydsk: Die bezau\-%
berte Welt.} % ***)
hvori
han bekæmper Troen paa onde Aanders Indvirkning, Djævlebesæt\-%
telser, Hexe o.\ desl. Stolende paa Fornuften efterviser han den
Uvidenhed, hvorfra de mange urimelige Beretninger om Djævelen
og hans Aabenbarelser i Hedenskab, Jødedom, Muhammedanisme
% --- p. 25 ---
\opage{25}og Christendom maatte antages at have deres Oprindelse. Det
Nye Testamentes Fortællinger om de Dæmoniske forklarer han
paa naturlig Maade og antager, at Jesus selv kun har villet refor\-%
mere det Moralske i Troen uden at rokke ved indgroede Fore\-%
stillinger. Om Fristelseshistorien udtaler han sig, idet han gaaer
ud fra Syndefaldet, blandt Andet saaledes: „Denn vor fest und
ungezweiffelt anzunehmen, dass ein geschaffener Geist, und der von
Gott verworfen ist, auff des Menschen Seele oder Leib wircken
könne, davon ist schon hier vornen angezeiget, wie viel daran lieget.
Hat der Satan eusserlich solch Gespräch mit dem Menschen ge\-%
führet, da die Welt erst den Anfang hatte, wie lieset man das
nicht mehr? Selbst nach Verlauff von viertausend Jahren in der be\-%
rühmten Unterredung mit dem andern Adam, da der Teufel zwar
ausdrücklich genennet wird, ist das nicht geschehen. Denn obschon
beydes Lucas und Mattheus das umständlich erzehlen, so muss
man sich dennoch wol hüten, dass man nicht alles nach dem Buch\-%
staben verstehe. Oder man müste zugleich glauben, dass Moses
und Elias beyde eben so wol auff dem Berge persönlich gewesen
und sich mit dem Seligmacher unterredet haben, wie solches ein
andermahl, ja dreymahl wird erzehlet. Matth.\ XVII, 3. Marc.\ IX,
4. Luc.\ IX, 30, 31. Gleichwol war Moses nicht leiblich wie Elias
in den Himmel auffgenommen worden, sondern gestorben und gar
gewiss begraben, denn Gott hat es selbst gethan. Deut.\ XXXIV, 6.
War er nun denn von den Todten auferstanden und leibhaftig gen
Himmel gefahren: Wunder, dass die Schrifft ein so grosses Wunder
nicht meldet. Wer denn von solchen, als im Gesichte geschehen
ist, so redet, als etwas das nach dem Buchstaben verstanden wird,
der weiss, wie Petrus, auch nicht was er saget“.\footnote{Anfr.\ Skr.\ S.\ 135--136 o.\ s.\ flg.} % *)
% sic: the note *) on p. 25 has no call mark anywhere in the text of the page; placed after the quotation it cites

Men seer vi heraf, at Fornuften ikke kan fordrage Djævelen
og derfor overalt, hvor han viser sig, falder ham i Flankerne, da
bliver det os tillige forstaaeligt, at den frie Fornuft, naar dens Øie
først er opladt, og Kritiken ret faaer begyndt, ikke kan indskrænke
sig til at forfølge Djævelen, men, orienterende sig paa alle Bibe\-%
% --- p. 26 ---
\opage{26}lens Enemærker, maa have Ret til at besinde sig paa, hvorledes
det kan gaae til, at Moses, der er bleven begraven af Gud selv,
og Elias, der er faren levende op til Himmelen, kunne aabenbare
sig paa eet og samme Bjerg; det er, naar vi see nærmere til, ikke
Djævelen alene, men Miraklerne, der er den frie Fornuft en For\-%
argelse. Man vil ligeledes indsee, at Fornuften, naar den først
har faaet fat paa Erkjendelsesprincipet, ikke kan være forskjellig
i de forskjellige Kirker, men maa hæve sig over de confessionelle
Begrændsninger. At den reformerte Kirke har en anden Troes\-%
bekjendelse end den lutherske, er vist, men hvo vil deraf udlede,
at Fornuften selv maa være en anden for Lutheranerne end for de
Reformerte? Den reformerte Præst blev afsat for sin „Bezauberte
Welt“, men hans Bog fandt ikke desto mindre Indgang hos alle
Fornuftige saavel i den lutherske (Semler og Thomasius) som i
den reformerte Kirke. I Danmark har ikke blot Holbergs Lune
ryddet betydelig op i Satans Papirer, men den med ham begyndte
Rationalisme har, om end med langsomme og forsigtige Skridt,
nærmet sig Alteret.

Paa Overgangen fra Orthodoxi til Rationalisme staaer den for
sin Tid fordomsfrie og respectable Tænker Christian Bastholm.
Bastholms Fornuft er i den Grad bleven færdig med Djævelen, at
denne end ikke udtrykkelig nævnes i Dogmatiken\footnote{Den Christl.\ Religions Hoved-Lærdomme. Kiøbhvn.\ 1783.}, % *)
% sic: the call in the text is *), but the note at the foot is marked **); there is no other note on p. 26
S.\ 265, hvor
Fristelseshistorien berøres. „Verdens Forfører“, hedder det, „ind\-%
stiller sig for ham“ (Christus); „mange og store Prøver sættes
paa hans Uskyldighed; men han tilfegter sig Seieren over dem
alle. Denne Fristelse var i adskillige Henseender nødvendig; han
maatte, førend han begyndte at omgaaes med Menneskene, lære at
kjende Mennesket, dets Skrøbeligheder og dets Anfægtninger, at
hans egen Erfarenhed kunde indgyde ham ømme Følelser for
Menneskenes Elendighed og deres Svagheder“ \ldots\ Man vilde
imidlertid tage mærkelig feil, hvis man heraf sluttede, at den Bast\-%
holmske Fornuft nu ogsaa var enig med sig selv om at bryde
Staven over alt Overnaturligt og forkaste Miraklerne. Idet For\-%
% --- p. 27 ---
\opage{27}nuften lister Erkjendelsesprincipet til sig, skeer det med en saadan
Varsomhed, for ikke at sige Behændighed, at den mindre Op\-%
mærksomme næsten Intet mærker. Om „Jesu Undergjerninger“
hedder det blandt Andet, S.\ 307 følgd., „at de oversteeg alle Na\-%
turens Kræfter og ophævede dens bekjendteste Love \ldots\ Vi be\-%
frygte aldeles ikke, at nogen engang skal bringe sin Kundskab i
Naturen saa vidt, at han skal beskæmme Jesu Gjerninger, at han
skulde kunne helbrede Syge blot ved en Befaling, blot ved en Haands
Paalæggelse, eller blot ved at røre ved ham, eller at han ved sin
blotte Villie skulde helbrede dem, som endog ere fraværende. Vi
befrygte aldeles ikke, at nogen engang skulde bringe sine Indsigter
i Naturen til den Høide, at han blot ved at true Bølgerne og
Stormvindene skulde standse dem midt i deres Raserie, eller, ved
at kalde ad Havets Fiske, skulde i samme Øieblik forsamle dem.
Dersom det var muligt engang ved Naturens Kræfter at gjøre de
samme Gjerninger, som Jesus Christus haver gjort, saa maatte
man dog allerede i vore Tider kunne finde nogle Spoer til dette
Haab, efter de stærke Skridt, man haver gjort, og den Høide, man
haver naaet i alle naturlige Videnskaber. Men hertil findes end
ikke det mindste Spoer“. „Men heri er jo intet Ufornuftigt; er
Christus virkelig den, han sagde sig at være, Aabenbaringen af
den guddommelige Sandhed, saa er det jo ganske fornuftigt, at
han har kunnet gjøre Gjerninger, som ellers overgaae alle menne\-%
skelige Kræfter“. „Denne Slutning“, hedder det, S.\ 318, „er saa
eenfoldig, at den Eenfoldigste kan fatte den. Derpaa grunder vor
Fornuft sin Roelighed, og sin Overbeviisning; mere behøver den
ikke at vide“. I Læren om „Apostlernes Undergjerninger“ hedder det,
S.\ 404, 405 o.\ følgd.: „Ved et Under forstaae vi en Forandring i
Naturen, som enten overstiger, eller er tvertimod Naturens Kræfter
og dens Love, i det mindste paa den Tid, da den skeer \ldots\ At
saadanne Under ere muelige, det kan ingen tvivle paa, som kjender
det uendelige Væsens Almagt. Naturens Love ere ikke i sig selv
nødvendige; de ere ikkun Udflydelser af Skaberens frie Villie, som
haver indrettet Naturen saaledes, som den er, haver skjenket
Skabningerne de Kræfter, som de besidde, og haver saaledes for\-%
% --- p. 28 ---
\opage{28}bundet Aarsagerne med deres Virkninger, som vi nu finde dem. Men
den, som haver givet en Lov, kan ogsaa ophæve den, naar viise
Hensigter fordre det“ \ldots\ „Haver Skaberen engang af frie Valg
givet Ilden den Kraft at lyse og at brænde, saa kan han og igjen
berøve den denne Kraft, eller hindre dens Virkning“ \ldots\ Dette tager
sig nu altsammen ret fornuftigt ud. Hvor fornuftig den hele
Aabenbaring, ifølge Bastholms Opfattelse, er bleven, kan blandt
Andet sees af Kapitlet om „Jesu Himmelfart“, S.\ 369--382.
„Efterat Christus havde omvandret 40 Dage med sine Disciple \ldots\ %
forlod han Jorden med sin synlige Nærværelse, for at tage Himmelen
i Besiddelse \ldots\ Det er den Orden, som Herrens Viisdom haver
fastsat, at efter Livets Trængsler følger Døden, efter Døden Op\-%
standelsen, og efter Opstandelsen en evig Besiddelse af den him\-%
melske Lyksalighed i Sjel og Legeme. Nu var Christus død, han
var opstanden; intet var altsaa naturligere, end at han nu skulde tage
Himmelens Saligheder i Besiddelse \ldots\ Fornuften finder altsaa intet
urimeligt i den Lærdom, at Jesus Christus \ldots\ blev optaget til
Himmelen, for at tage den Herlighed og Lyksalighed i Besiddelse,
hvilken tilkom ham som Gudmenneske \ldots\ Kan man her, ligesom
ved Christi Opstandelse, bevise, at det var umueligt, enten at Apost\-%
lerne ved denne Begivenhed kunde selv blive bedragne, eller at de
kunde ville bedrage os med et falsk Vidnesbyrd, saa er den Troe,
vi sætte til deres Efterretning herom, grundig og fornuftig“.
Christus har, ifølge Forfatterens Mening, anvendt fire Midler for
at overbevise Apostlene og Efterverdenen om, at intet Bedrag har
fundet Sted, at det, som her blev seet ikke har været en For\-%
blindelse af Sandser og Indbildning, men fuld Virkelighed. „Først
var dette et Stykke af Jesu Viisdom, at han ikke forsvandt iblandt
dem. Efter hans Opstandelse pleiede han undertiden gandske ufor\-%
modentlig at indstille sig iblandt dem, og igjen ligesaa ufor\-%
modentlig at forsvinde for deres Øine; maaskee han herved vilde
overbevise dem om, at det Legeme, han ved Opstandelsen havde
ført ud af Graven, havde andre og høiere Egenskaber, end det
Legeme, man nedlagde i Graven. Men nu, da han vilde forlade
Jorden, fordrede et andet Øiemeed et andet Forhold \ldots\ Vor For\-%
løser blev langsom og synlig optaget til Himmelen \ldots\ Det
% --- p. 29 ---
\opage{29}andet Stykke af Jesu Viisdom var dette, at han udvalgde Olie\-%
bjerget, for derfra at holde sin Himmelfart. Intet Sted kunde
være beqvemmere til dette Øiemeed. Dette Bjerg laae en Sabbats
Reise fra Jerusalem, de havde altsaa Leilighed nok til at see og
betragte Jesum i en frie og aaben Luft. Han talede med dem underveis,
gav dem sine Befalinger, igjentog sit Løfte om den Hellig Aand,
som han vilde sende dem, og velsignede dem \ldots\ Dette Bjerg
var og det høieste Bjerg omkring Jerusalem; her var altsaa en
frie Udsigt og en reenere Luft, som forskaffede dem en beqvem
Leilighed til gandske ubehindret at beskue saa vigtig en Tildragelse“.
Det tredie Stykke af Jesu Viisdom maa især forekomme os at have
et overveiende Præg af Fornuftighed. „Man spørger, hvorfor
Christus ikke her haver forsamlet alle sine Fjender, for ved saa
stort et Under paa een Gang at overbevise dem alle om deres
Vantroe og hans guddommelige Sendelse? Jeg kunde svare meget
paa dette Spørgsmaal; jeg kunde svare, at de vel lode sig lige\-%
saa lidet overbevise ved dette Under, som ved alle de øvrige \ldots\ %
Jøderne tilskreve Djevelen hans andre Undergjerninger; de kunde
altsaa sige det samme om denne. Naar et Menneske er i den
Grad forhærdet, saa nytte Beviser ikke mere \ldots\ Jeg seer altsaa
ingen Aarsag, hvorfor Christus skulde opfare i sine Fienders Paa\-%
syn; men jeg seer en vigtig Aarsag, hvorfor han ikke skulde gjøre
det. Dersom Christus havde offentlig bekjendtgjort, paa hvad Tid
og hvad Sted han vilde opfare til Himmelen, saa havde han unæg\-%
telig draget alle Jerusalems Indbyggere efter sig; thi hvo føler
ikke Nysgierrighedens Magt, og hvo skulde ikke gaae en liden
Vei for at være Tilskuer af saa overordentlig en Begivenhed?
Men hvilken Larm, hvilken Tummel og Uorden maatte der ikke
nødvendig have været i saa talrig en Sværm, hvor den ene trænger
den anden, hvor enhver vil være den nærmeste, og den ene fordriver
den anden, ligesom paa et oprørt Hav den ene Bølge den anden“.
Vistnok Noget af det Fornuftigste, der kan siges i en videnskabe\-%
lig Dogmatik.

Hos Bastholm ligger endnu Fornuftprincipet, om vi saa maae
sige, indsvøbt i det Supranaturale og sparker i Svøbet; men i
% --- p. 30 ---
\opage{30}den tydske Literatur er Fornuften allerede tidligere blevet sig sin
Overlegenhed bevidst. Den Wolfiske Philosophi har i den Grad
gjennemtrængt den Wertheimske Bibeloversættelse\footnote{Erster Theil, 1735.}, % *)
at Guds Aand
ikke kan svæve over Vandene, men „en hæftig Vind“ maa begynde
at blæse; Begreberne defineres saa skarpt, at der til Forstaaelse
af 3die Moseb.\ XVIII, 7 endog for Tydeligheds Skyld maa tilføies,
hvad en Moder er, nemlig „en Kone, som avler Børn i Selskab
med en Mand“. Boghandleren og Forfatteren Nicolai i Berlin
udgav fra 1765 til 1805 „Allgemeine“ og „Neuallgemeine deutsche
Bibliothek“, hvori Fornuften blev ganske populær. Theologerne
selv arbeidede med paa Tydsklands Oplysning. Den Cölnske Provst
Wilh.\ Albr.\ Teller gav et offenligt Svar paa et Sendebrev fra
berlinsk-jødiske Huusfædre, hvori han erklærede sig villig til at
anerkjende dem som Christne paa deres blotte Deisme. Større
Liberalitet kunde man da ikke vel forlange; det er jo, som om vi
hørte Fornuftens egen Røst. „I Aaret 1833 f.\ Ex.\ kaldte den
evangeliske Generalsuperintendent Röhr's „Predigerbibliothek“ alle
christne Kirkers Lære om den treenige Gud et „\emph{antichristeligt
Dogma}“ (1833, H.\ 1. S.\ 11), og Augustin formedelst hans Lære
om Synden og Naaden, og fordi han ikke var forbleven uomvendt,
En, „som vanhelligede Evangeliet, og et udsvævende Menneske,
som var bleven en af de Hellige“ (sammested S.\ 13), Pelagius
derimod (S.\ 15) „Fornuftens ærværdige Forsvarer mod Ufornuften,
en Mand, som allerede paa sin Tid havde den Tilfredsstillelse, at
see de Viseste og Bedste staae paa sin Side“, (prisende den evan\-%
geliske Kirke lykkelig (S.\ 16) naar Luther ikke blot havde refor\-%
% "samt erklærede": a risen space (spacing material inked) between the two words, not a letter
meert tilbage til sin Ordensfader Augustin), samt erklærede (S.\ 48):
„Supranaturalismen fører bort fra Gudstjenesten til Afgudstjene\-%
sten“ o.\ s.\ v.; ja, endnu i det samme Aar 1833, fremstillede den
ovennævnte høie Geistlige, og det ikke blot \emph{privatim}, men al\-%
deles offentlig og med \emph{kirkelig Lovskraft} hele Christendommens
Fundament, Ordet om Korset, Læren om Christi Død, „\emph{en til
Menneskenes Bedste} („til menneskelig Ondskabs Udsoning“)
% --- p. 31 ---
\opage{31}udholdt Forløsnings- og Forsoningsdød“, som „en Følge af det rene
Evangeliums totale Forvanskning“, som „\emph{uvidende og spids\-%
findige Kirkelæreres Digt}“. Ja, man gik endnu videre,
--- fra Halle forkyndte man alle Dogmers Fald (Halliche Allg.\ %
Lit.\ Zeitung; März 1833 Ergänz.\ Bl.\ S.\ 234, udtaler bestemt det
Haab at „\emph{naar alle christelige Dogmer seent eller tid\-%
lig falde}“, ville alle Mennesker endelig forene sig om en christelig
Moral“).\footnote{H.\ E.\ F.\ Guerike. Haandbog i Kirkehistorien. Overs.\ af Sverdrup.
Christiania. 1842. Anden Deel. S.\ 338 Anm.}. % *)
% sic: full stop printed both before and after the call: "Moral“).*)."

Den vulgære Rationalisme var imidlertid for overfladisk til
at kunne tilfredsstille dybsindige Aander; dertil udkrævedes en
speculativ Fornuft, som var istand til at gjennemskue det Over\-%
naturliges Identitet med det Naturlige. Dogmatikere som Daub
og Marheinecke saae ned paa Rationalisterne, men op til Schelling
og Hegel. Her fandt man en Fornuftudvikling saa dyb og spæn\-%
dig, at den rummede hele Christendommen. Da Forestillingen
var opløst i Begreb, blev Theologien med Eet saa orthodox, at
Djævelen selv gik med i Begrebet. Da Fornuften, ved at gjøre
sit Vidensprincip til Aabenbaringsprincip, just stod i Begreb med
at faae hele Orthodoxiens positive Indhold tillivs, blev den saa
speculativ opsvulmet, at den var nær ved at gaae fra Forstanden.
Det kunde virkelig have seet betænkeligt ud med dens dyrebare
Liv, hvis man ikke i rette Tid var kommet den Svulmende til
Hjælp ved et radicalt og rensende Middel; dette Middel var en
afførende Kritik, der navnlig blev anvendt af Dr.\ Strauss.

Hvor meget man end kan beklage, at Strauss har været
blind for Religionens guddommelige Væsen, saa maa man dog paa
Videnskabens Vegne være ham Tak skyldig, at hans kritiske Be\-%
handling saa vel af Jesu Liv som af den kirkelige Dogmatik, om
der end i det Enkelte kan være Adskilligt at indvende, har viist,
at den religiøse Tro maa gaae til Grunde, naar den skal retfær\-%
diggjøres i Viden, d.\ v.\ s.\ naar Religionen skal opfattes og be\-%
% the quotation opened here continues onto p. 32
grundes videnskabeligt. „Strauss' Talent bestaaer i at afdække
% --- p. 32 ---
\opage{32}Forvirringen og tilintetgjøre Illusionerne; men hans Kritik er kun
% letterspacing in the print (opløsende, negativt) -> \emph; a speck before "opløsende" is dirt
\emph{opløsende}, Resultatet forbliver \emph{negativt}“. „Hans Dogmatik
er ingen Dogmatik, men kun en Kritik af de enkelte Dogmer, et
Repertorium af dogmatiske Forestillinger. Trods al Reenlighed i
Stoffets ydre Ordning og Begrændsning, og trods Forstandsretnin\-%
gens Sikkerhed, er dog en uhyre Mangel kjendelig, og den nihi\-%
listiske Baggrund frembringer en Følelse af Trøstesløshed og
Tomhed“\footnote{Dr.\ Karl Schwarz. Den nyeste Theologies Historie. Dansk Over\-%
sættelse. Kjøbenhavn. 1879. S.\ 779.}. % *)
Denne Opfattelse maa vistnok billiges af Alle, der
af Videnskaben endnu forvente en ny Theologi, men det er netop
denne Forventning, som den Straussiske Kritik udtrykkelig er lagt
an paa at tilintetgjøre. At han i denne Henseende har løst sin
Opgave, er aldeles utvivlsomt. Man kan omdanne Rationalismen
saa meget man vil, det rationelle Princip vil aldrig blive religiøst.
K.\ Schwarz anpriser Hases rationelle Theologi. „Det Alt beher\-%
skende Princip i hans Dogmatik“, siger han\footnote{Anfr.\ Skr.\ S.\ 413.}, % **)
„var jo Menne\-%
skets \emph{relative Frihed}, den paa denne grundede Kjærlighed
til Guds Rige. Den Treenighed, han udviklede, var ogsaa en ganske
anden end Schellings, der forvandlede sig til en fysisk Verdens\-%
begivenhed, en theogonisk Proces, hvorimod den hos ham var ethisk
og blev ført tilbage til det praktisk-bibelske Indhold. Overhovedet
hævdede han med fuld Bevidsthed sit theologiske System den
ethiske Charakter, og grundede det paa Menneskets sædelig-reli\-%
giøse Frihed i hans Selvstændighed for Gud, samt paa hans Kjær\-%
lighed og uendelige Stræben efter Samfund med Gud, som paa den
% doubtful: a tiny dot after "Princip" (smaller than a stop, at 600 ppi) taken as dirt; rejected: comma (text layer)
fasteste Grund. Dette ethiske Princip er jo Rationalismens ufor\-%
gængelige Sandhedskjerne“.

Men det ethiske Princip er jo ikke noget særligt videnskabe\-%
ligt Princip. Fornuften kan postulere Frihedens Realitet og an\-%
befale Menneskene i det praktiske Liv at antage den og troe paa
den; men det forholder sig ikke med Friheden som med Natur\-%
tingene. Alt hvad Videnskaben skal kunne bevise, maa den be\-%
vise i Kraft enten af absolut eller betinget Nødvendighed; men
% --- p. 33 ---
\opage{33}naar Friheden selv begrundes paa Nødvendighed, saa kan jo den
videnskabelige Fornuft med fuld Ret betvivle dens Realitet. Endnu
mindre kan Kjærlighed til Gud begrunde nogen videnskabelig For\-%
visning om Guds Kjærlighed, Guds Personlighed, saa lidt som en
personlig Trang til et evigt Liv kan afgive videnskabeligt Beviis
for Sjælens Udødelighed. At Fornuften med sit Princip virkelig
skulde begrunde en Treenighedslære, som videnskabelig førte til\-%
bage til det „praktisk-bibelske Indhold“, er en Paastand, der for\-%
raader en aldeles uklar Sammenblanding af Troens og Videns
Charakteer, og stadig forvexler religiøse Forestillinger med viden\-%
skabelige Kjendsgjerninger, Aabenbaringstanker med almene For\-%
nufttanker.

Religionens Gud er, hvorledes man end vender og dreier det,
som Frihedens og Kjærlighedens Gud et overnaturligt Væsen.
Kan det nu med Sandhed siges, at Fornuften i al Videnskabelig\-%
hed ikke har noget Blik for det Overnaturlige, saa har den jo
heller intet Blik for Frihedens, Personlighedens og Kjærlighedens
Gud. Fornuften maa med Forbigaaelse, ja for saa vidt den skal
gjøre sit Princip gjældende, med Fornægtelse af det Mirakuløse
i Skabelsen saa vel som i Incarnationen og i Inspirationen, holde
sig strængt til Naturen og dens Love. Hvor vil den „rationelle
Theologi“ da hen med sin Gudskjærlighed og sit ophøiede Friheds\-%
ideal? Den maa søge sin Guddom i Naturen og binde den til
Naturlovene, den maa opgive Friheden og forgude Nødvendigheden.
--- Den, der vil hylde Comtes Religion, bliver dog nødt til enten paa
Maa og Faa at gaae videre end Fornuften, eller vende om og
bøie sig for Littrés Kritik.

% print: a short centred rule separates Section II from the head; head numeral in heavier type
\begin{center}III.\end{center}

At Religionen ikke kan tilegne sig Erkjendelsesprincipet i
hele dets Omfang og derved drage Videnskaben ind under sig, er
% --- p. 34 ---
% sic: "uu" (turned n) for "nu", as printed
\opage{34}viist i det Foregaaende; Spørgsmaalet bliver uu, om da Viden\-%
skaben med sit Fornuftprincip kan magte Religionen eller, om den
maaskee ved at forfølge sit eget Maal og opdage en væsenlig
Grændse vil komme til at indsee Nødvendigheden af at indrømme
Religionen et eiendommeligt Troesprincip. For at see, hvorledes
det hermed forholder sig, kaste vi et Blik paa Videnskabens Veie.
Den har kæmpet og maa endnu med sit Princip bestandig kæmpe
sig frem igjennem en dobbelt Række af Modsigelser, den ene
uundværlig, den anden i høi Grad hindrende og besværlig. Til
de uundværlige Modsigelser henregne vi alle dem, som have deres
Udspring fra sand Interesse for Videnskabens Fremme, Indsigelser,
ved hvis Besvarelse begaaede Feil hæves og Methoden sikkres;
til de hemmende, besværlige og belæmrende maae vi derimod hen\-%
regne de Modsigelser og Indsigelser, som hidrøre fra Ukyndige,
der af en og anden Interesse gjerne ville standse Videnskaben.
Det er, for Principets Skyld, især de Sidste, hvorpaa vi maae hen\-%
lede Opmærksomheden.

Nikolaus Copernicus, som allerede 1530 havde sit nye Ver\-%
denssystem\footnote{De Revolutionibus orbium coelestium, Norimbergæ MDXLIII.} % *)
færdigt, lod det dog først trykke 1543 og berørte
det første Exemplar med døende Haand. At det ikke saa meget
har været Frygt for de Ptolemæiske Astronomers som meget mere
for Theologernes og Præsternes Polemik, der har holdt ham til\-%
bage, kan sees af hans Dedication til Pave Poul III, idet han
% Greek fount in the print; Latin below in the same roman as the Danish
forudseer, at tomme Ordgydere (ματαιόλογοι), blottede for al ma\-%
thematisk Viden, ville anmasse sig en Dom over Værket ved med
Villie at fordreie et eller andet Sted i den hellige Skrift (propter
aliqvem locum scripturæ male ad suum propositum detortum);
denne Anelse gik da ogsaa rigelig nok i Opfyldelse. Tycho Brahes
Observationer, Keplers Love, Galileis Opdagelse af Jupiters Maaner
% sic: "1689" as printed
og endelig Newtons mathematiske Principer (1689) betegne stor\-%
artede Fremskridt i en Periode, hvor ellers kirkelig Orthodoxi
holdt Videnskaberne lænkebundne. Vistnok maatte Theologerne
være forbittrede paa disse ugudelige Theorier, thi endskjøndt de
% --- p. 35 ---
\opage{35}ikke rigtig forstode dem, var dog saa meget klart, at de maatte
tjene til at undergrave Bibelens Autoritet; men hvad skulde de
vel gjøre? Himmelhvælvingen brast, nye Solsystemer viste sig;
Jorden svandt ind paa en for Christus og Christne høit upassende
Maade. Tilsidst, ved Overgangen fra attende til nittende Aar\-%
hundrede, trængte man Skaberen tilbage, lod Kloderne danne sig
af en Urtaage og holdt sig ved alle Dannelser til Materie, Kræfter
og Love.

Vistnok var ikke hermed Alt afgjort. Kunde Theologen end
ikke følge Astronomen paa de steile Veie mellem Stjernerne, saa
kunde han dog trøste sig med, at det menige Folk heller ikke
kunde det. Vilde det aldrig blive Folket ret indlysende, hvad
Astronomernes videnskabelige Udskeielser havde at betyde, saa
kunde man jo ignorere dem, naar man blot passede paa at holde
Almanakken reen. Det gjaldt da om at møde Videnskaben, ikke
i det umaalelige Himmelrum, men paa den faste Jord, og her at
forhindre dens Afvigelse fra Aabenbaringens Ord i Skriften. Men
ogsaa her har Videnskaben taget Magten fra Theologien. Vistnok
er Geologien, Læren om Jordens trinvise Udvikling, endnu saa
ung, at dens Skikkelse idelig kan forandres; men Saameget er
allerede klart, at den ikke retter sig efter den bibelske Skabelses\-%
lære, men efter Kjendsgjerningerne i Naturen. De Modsigelser,
den har at kæmpe med, falde allesammen inden for Videnskaben
selv; Indsigelser fra et overnaturligt Standpunkt ændser den slet
ikke. Det er vistnok kun sjældent, at Theologerne for Alvor ind\-%
lade sig med Geologien, men naar det leilighedsviis skeer, for\-%
raade de gjerne en kolossal Misforstaaelse. Saaledes taler f.\ Ex.
Theologen Tholuck\footnote{Vermischte Schriften. Zweiter Theil. Hamburg 1839. II. Was ist
das Resultat der Wissenschaft in Bezug auf die Urwelt? S.\ 148
o.\ s.\ frmdl.} % *)
stærke Ord især i Anledning af et Par andre
% German set in the same roman as the Danish (no change of fount), so no \textit
Theologers anerkjendende Udtalelser om Geologien. „Die Er\-%
fahrungswissenschaften aller Art, heisst es\ldots\ (Bretschneider i s.
Sendschreiben), haben fühlbarer und störender in den Bestand des
% --- p. 36 ---
\opage{36}alten theologischen Lehrsystems eingegriffen, als die spekulative
Philosophie\ldots\ Die Geologie will die mosaische Schöpfung nicht
mehr mit den Revolutionen, die unser Erdball erlitten hat, ein\-%
% doubtful: after "unbekümmert" a tiny baseline dot (a third the size of a stop, 600 ppi), read as a comma whose tail did not print; rejected: full stop (text layer)
stimmig finden. Sie lehrt, unbekümmert, wie der Theologe dabei
% sic: "grosze" as printed
zurechtkomme, dass die Erde mehrere grosze Bildungsepochen von
unbestimmbarer aber langer Zeit durchgegangen ist, und dass die
ersten Schöpfungen auf ihr wieder untergegangen sind“. „Die er\-%
habene Wissenschaft der Astronomie war es\ldots, welche in die Be\-%
griffe des Alterthums von Himmel, Erde, Hölle, Auferstehung,
Gericht, Ende der Welt, die noch zur Zeit der Reformation un\-%
verändert waren, auflösend eingriff --- die Naturforscher und Reise\-%
beschreiber schilderten die Verschiedenheit der Racen an Gestalt,
Farbe und geistigen Kräften, die durch die Vermischung der
% sic: "entshehenden" as printed
Racen entshehenden Spielarten, und wiesen die grossen und blei\-%
benden Unterschiede unter ihnen nach, indem sie zeigten, dass
diese Differenzen nicht auf Rechnung des Klimas und der Nah\-%
rung, sondern auf \emph{Verschiedenheit der Grundabstam\-%
mung} sich gründen müssen. \emph{Blumenbach} sammelte die
Schädel in allen Welttheilen und brachte die Ansicht hiervon in
ein System. In welche Verlegenheit gerieth nun der Theologe?
Wenn es nun nicht mehr \emph{einen} Adam für alle Menschen, son\-%
dern einen Adam für die Kaukasier, einen andern für die Neger,
einen dritten für die Amerikaner, einen vierten für die Malayen,
einen fünften für die Mongolen u.\ s.\ w.\ gegeben hat; wo blieb
% sic: "mit den einen" as printed
nun die Dogmatik mit den \emph{einen} Adam der Bibel, mit der Lehre
vom Sündenfall und von der durch Adam auf alle Menschen ge\-%
brachten Schuld, wo nun mit der ganzen Lehre von der Erbsünde
als Folge des Falles und einer von Adam aus durch Zeugung an
alle Menschen gekommener Schwäche? Und ging diese verloren,
wie stand nun die Nothwendigkeit der stellvertretenden Genug\-%
thuung Christi, des zweiten Adams, um die Schuld des ersten
Adams aufzuheben, zu erweisen? Wo blieb nun der Grund der
Verdammniss der Heiden, die nicht von Adam abstammen?“ ---
% sic: "Tholuch" as printed
Over denne Bretschneiders uforbeholdne Udtalelse er Tholuch
yderlig indigneret. „So lassen sich“ --- udbryder han --- „die
% --- p. 37 ---
\opage{37}Diener der christlichen Kirche im neunzehnten Jahrhundert nach
Entstehung derselben vernehmen. Sieht man, nachdem eines \emph{Juli\-%
anus}, \emph{Porphyrius}, \emph{Voltaires} Stimme verklungen ist, die Pfleger
\emph{des christlichen Heiligthums} mit solchen Reden auftreten,
wer muss nicht ausrufen: Herr! Bewahre deine Kirche vor ihren
Freunden und Pflegern, mit ihren Feinden hat es keine Noth!“

Idet nu Tholuck afmønstrer de vigtigste Geologer fra Cartesius
indtil sin egen Tidsalder og seer Palæontologien efter i Sømmene,
glemmer han naturligviis ikke at anføre en Cuvier, en Leopold
von Buch, en Humboldt o.\ fl., men kommer netop derved til fuld
Forvisning om sit eget theologiske Resultat, at det Hele Intet har
at betyde. Theologien er en sand Videnskab; men Geologien med
sine mange Hypotheser er egenlig ikke nogen Videnskab. „Der
christliche Theologe kann in der That, so lange die geologische
Wissenschaft keinen anderen Standpunkt als den gegenwärtigen
einnimmt, in Bezug auf alle jene Hypothesen nur das stolze Wort
gebrauchen, welches jene Königin von Schweden, Christine, sprach,
da sie ihre Krone niederlegte: \emph{non mi bisogna e non mi
basta --- ich bedarf ihrer nicht und sie ist mir nicht
genug}“. Naar Tholuck, for at forklare Skabelsen, beraaber sig
paa Guds Aand, der svævede over Vandene, og for Øvrigt hen\-%
viser til Poesien\footnote{Anfr.\ Skr.\ S.\ 172.}, % *)
da er kun det herved mærkeligt, at han ikke
selv kan mærke, at han er kommen langt ud over Videns og Vi\-%
denskabens Grændse.

Hvad det Palæontologiske angaaer, har Darwinshypothesen
i de sidste Decennier ryddet saaledes op i Botanik og Zoologi, at
det maaskee om mange dygtige Fagmænd kan være vanskeligt at
sige, paa hvilken Side de egenlig staae; kun saa meget er klart,
at de neppe kunne staae paa Læren om en overnaturlig Skabelse.
E.\ Hæckel, en af Darwinismens meest fremragende Tilhængere,
har ved Lærdom, Skarpsindighed og systematiserende Talent givet
Anviisning paa at forklare Livets Oprindelse af physiske Aarsager,
Betingelser og Love, uden nogensomhelst Indblanding af en fore\-%
% --- p. 38 ---
\opage{38}given Teleologi, og ved at paavise de nødvendige Overgangsled
forfulgt Dyrenes Afstamning fra Ormen til Mennesket i een con\-%
tinueerlig Række. Hans Forklaring af Livet skulle vi komme til\-%
bage til; hvad hans Systematik angaaer, ville Zoologerne i Frem\-%
tiden vistnok ikke blot faae Meget at anerkjende, men ogsaa
Meget at rette. En sammenhængende, fremadskridende Udvikling
af de organiske Naturformer kan jo meget vel have sin Gyldighed,
uden, at det derfor strax lykkes Videnskaben at paavise de særlige
Overgangsled; at digte dem giver kun subjectiv, ikke objectiv Til\-%
fredsstillelse. Ved en Behandling af det Spørgsmaal: „\emph{Ned\-%
% sic: no closing quotation mark after "Krybdyrene?" in the print
stamme Fuglene fra Krybdyrene}?\footnote{Skildringer af Dyrelivet i Fortid og Nutid. Kjøbenhavn 1880.} % *)
har Dr.\ Chr.\ Lütken
leveret et Exempel paa, hvorledes den grundige, sindige Forsker
fatter sine Opgaver. Efter at have fulgt Fuglenes Spor saa langt
som muligt tilbage i Jordens Udviklingshistorie, gjennem Kridt-
og Juraperioden ned til Trias, ender han med det Resultat, at
Fuglenes Tilnærmelse til Krybdyrene er for Lidet eller Intet at
regne. „Maaske, siger han S.\ 41, skulde vi søge Ligheden fra
den modsatte Side, fra \emph{Krybdyrene} over imod Fuglene? Det
er i saa Fald til Reptilklassens Fortid, at vi maa vende os, og vi
ville ikke længe være i Tvivl om, hvad det er for Former, som
særlig ville gjøre Krav paa vor Opmærksomhed; det maa være
de flyvende \emph{Krybdyr}!“ Men betragte vi nu \emph{den langsnudede
Flyveøgle} (Pterodactylus longirostris), saa viser det sig dog, at
Haabet skuffer. „Først Cuviers overlegne Skarpsindighed og ge\-%
niale Blik“, hedder det S.\ 45, „satte ham i Stand til efter de
ufuldkomne Tegninger, som vare ham tilgængelige, at udvikle, at
% sic: the inner quotation opens with a raised » (wrong sort) and closes with “, as printed
det var \emph{et flyvende Krybdyr}, en »Flaggermus blandt Øg\-%
lerne“, men ikke nogen Flaggermus i sædvanlig Betydning, over\-%
hovedet intet Pattedyr, saa lidt som nogen Fugl“. Da det nu i
det Hele ikke vil gaae med Flyve-Øglerne, har man forsøgt det
med Kæmpe-Øglerne (Dinosaurierne). „Saa nær“, siger vor For\-%
fatter, ifølge forudgaaende Oplysninger, S.\ 99: „vare da Dinosau\-%
rierne langt om længe komne Fuglene i Palæontologernes Opfat\-%
% --- p. 39 ---
\opage{39}telse, at det maatte ansees for overmaade tvivlsomt, om man var
i Stand til at skjelne deres Fodspor fra samtidige Kæmpe-Struds\-%
fugles! Da fremkom omtrent samtidig, for en halv Snes Aar
siden, Prof.\ \emph{Cope} i Nordamerika og Prof.\ Huxley i England med
Opdagelser og derpaa støttede Opfattelser, ifølge hvilke Dinosau\-%
rierne og Strudsene skulde vise saa store Tilnærmelser og Over\-%
eensstemmelser, navnlig i \emph{Baglemmernes Bygning}, at Di\-%
nosaurierne maatte siges, \emph{i det mindste i denne Henseende},
at staa midt imellem Fugle og sædvanlige Krybdyr (Krokodiler
og Nutids-Øgler).“ Det er, navnlig med et Blik paa \emph{Baglem\-%
mernes Omdannelse}, at Forfatteren advarer imod at blive
for tidlig færdig med Afstamningslæren. „Maaske“, siger han
S.\ 111, „er Seiren forberedt, ogsaa paa dette Omraade, og kan
komme førend vi vente den. Men at sende Seirsbudskab ud i
Verden, er ikke det samme som virkelig at have seiret“. Det er
med særligt Hensyn paa det videnskabelige Erkjendelsesprincip, at
vi anføre dette. Man gaaer ikke tilbage til mirakuløse Forud\-%
sætninger, fordi man i Videnskaben modarbeider overilede Theo\-%
% doubtful: "Forfatter" -- the second t has lost its crossbar and reads like l; rejected: "Forfatler"
% sic: the quotation is not reopened after "endog gjærne," but closes at "Vanskeligheder“" -- one “ without „
rier. „Vi indrømme“, erklærer vor Forfatter S.\ 5, endog gjærne,
at skulle vi træffe Valget mellem eet af to: enten at antage, at
alle Arter ere skabte pr.\ Mirakel (være sig gjennem en Række,
periodisk sig gjentagende, samlede Fornyelser af den hele organiske
Natur, eller ved en Uendelighed af isolerede Skabelser), \emph{eller} at
antage, at Nutidens Arter have udviklet sig af den foregaaende
Periodes, dennes igjen af den næstforegaaendes osv., da er Valget os
% sic: comma (tail clear at 300 ppi) before capital "Vi"
ikke vanskeligt, Vi foretrække naturligvis det \emph{mulige} (mere
eller mindre sandsynlige eller usandsynlige) for det videnskabelig
talt absurde, selv om vi ikke skjønne, at man derved er kommet
ud over alle Vanskeligheder“.

Vanskelighed af særegen Art frembyder Biologien. Det er
nemlig ikke alle Physiologer, der, saaledes som Hæckel, have Mod
til at fornægte al væsenlig Forskjel mellem Livløst og Levende,
en Forskjel, der skal kunne reduceres til visse Ændringer af
Molekuler og physiske Kræfter. Det kan udtrykkelig godtgjøres,
at Kræfterne i Mechanik, chemisk Physik og Chemi virke i Med\-%
% --- p. 40 ---
\opage{40}før af Egenskaber i Stoffet, der ligefrem svare til Inertiens Lov
og desaarsag ikke forraade Spor af Liv. At Livsyttringer i de
laveste, ufuldkomneste Organismer, Monerer f.\ Ex., synes at være
blot physiske Virkninger af physiske Kræfter, er et let forklarligt
Bedrag, der forraader sig selv, naar man fra et høiere Livstrin
gaaer tilbage til lavere. I Hæckels „Generelle Morphologie der
Organismen“ (S.\ 234, 1ste Deel) hedder det ganske naivt, at
Dyrets Sjæl kun beroer paa, at Forestillinger vækkes i Ganglie\-%
cellernes Æggehvidemolekuler, d.\ v.\ s., naar disse Molekuler selv
begynde at forestille; den Forestilling, der vækkes ved de \emph{cen\-%
tripetale} Traade, er \emph{Fornemmelse}, altsaa: Molekulerne for\-%
nemme; den Forestilling, der vækkes i de \emph{centrifugale} Traade,
er \emph{Villie}, altsaa: Molekulerne ville; den Forestilling, der ud\-%
vikles under meget complicerede Vexelvirkninger af centripetale
og centrifugale Traade, er \emph{Tanke}. Det er da ikke Ganglier og
Hjerne, der blot afgive nødvendige Betingelser for Tænkningen,
men Gangliecellernes Æggehvidemolekuler ere selv tænkende:
Hjernen er Tænkningens egenlige Subject. Men hvor sandt det
end er, at Videnskaben hverken kan paavise, hvorfra det organiske
Liv egenlig kommer, ei heller tydeliggjøre, hvorledes den punc\-%
tuelle Vexelvirkning med Materien foregaaer: saa er det dog ind\-%
lysende, og kan videnskabelig godtgjøres, at et Indbegreb af ele\-%
mentære Kræfters Virksomhed ikke ved nogensomhelst Ændring
eller Omdannelse lader sig forvandle til Livsvirksomhed, at enhver
levende Organisme netop maa beroe paa concret individuel Eenhed,
Forskjelseenhed, oprindelig Eenhed af modsatte Functioner, Stof\-%
functioner og Selvhedsfunctioner\footnote{Almindelig Videnskabslære. Kjøbenhavn 1880. S.\ 172 flgd.}. % *)
Er det end, ifølge den menne\-%
skelige Videns væsenlige Begrændsning, ikke muligt at naae videre,
saa kunne vi dog, ved at besinde os paa det almindelige Forhold
af Subjectivt og Objectivt, der gaaer igjennem den hele Natur,
forvisse os om, at Livets Fremkomst og Udvikling beroer paa
særegne Anlæg, der ingenlunde bryde den videnskabelige Sammen\-%
hæng ved at kræve Mirakler.

% --- p. 41 ---
\opage{41}„\emph{Fra høiere Instans til lavere ei bring din Sag}“,
disse Ord af en engelsk Physiolog, der ikke vil gjøre Physiologien
umiddelbart afhængig af Chemien, finde især Anvendelse paa For\-%
holdet mellem Biologien og de elementære Videnskaber. Skjøndt
Physik og Chemi ere saa vel begrundede Videnskaber, at Ingen
kan betvivle deres Gyldighed, saa er det dog unægteligt, at Grund\-%
laget, hvorpaa de hvile, nemlig Atomer og Kræfter, ingenlunde
kan siges at være enten ved sig selv indlysende eller i og for sig
tilstrækkeligt. Et System af Verdenslove maa nødvendig forud\-%
sættes, naar en Videnskab skal være mulig, og det kan ikke, som
Kant vil, være vor egen Bevidsthed, hvorfra Verdenslovene udgaae;
men tænke vi al Bevidsthed borte, saa er der lige saa lidt Love
som Atomer og Kræfter. Eller skulde maaskee det Objective
kunne være, ikke blot forud for vor endelige Subjectivitet, thi det
forstaaer sig, men forud for al Subjectivitet, altsaa før det blev
objectiveret? Det er umuligt. Lovene kunne ikke være uden
Tanker, efterdi de have Almeenbestemmelserne i sig; men de
kunne heller ikke tænke sig selv som Kræfter i Materien: her
maa forudsættes en mægtig tænkende Selveenhed, en evig Sub\-%
jectivitet, hvori den objective Verdens Tilværelse er begrundet.
Lad os kun tilstaae, at der i det oprindelige Forhold mellem Sub\-%
jectivt og Objectivt, mellem Indre og Ydre, Ideelt og Reelt, mel\-%
lem Tingenes Begreber og Tingenes Existens er noget for vor be\-%
grændsede Fornuft Uigjennemtrængeligt, saa viser sig dog det al\-%
mindelige Grundforhold selv desto klarere, jo mere vi tænke over
Vexelbestemmelsen af Ting og Love, Virkelighed og Mulighed,
Totalitet og Anlæg. Maa Solsystemet forudsætte et Anlæg, hvori
det Objective oprindelig er bestemt ved det Subjective, maa den
geologiske Udvikling beroe paa et særligt i Solsystemet præfor\-%
meret Anlæg, da maae de levende Organismer ogsaa vise tilbage
til Anlæg, der ligge i Subjectivitetens Væsen. Saa meget staaer
fast, at det i Organismerne Subjective ikke kan afledes af noget\-%
somhelst Objectivt, men viser tilbage til Subjectivitetens evige
Væsen; men hvorledes den uendelige, Lovenes og Tingenes objec\-%
tiverende Subjectivitet indformer de subjective Anlæg i Stoffet, er
% --- p. 42 ---
\opage{42}og bliver en uigjennemtrængelig Gaade: vi staae ved en uover\-%
stigelig Grændse.

Kan det nu ikke videnskabelig godtgjøres, hvorledes Livet
kommer frem i Materien og virker; maa her, hvad dets Oprindelse
angaaer, nødvendig blive en Dunkelhed tilbage; med hvad Ret
vil man da kunne forlange, at Videnskaben nøiagtig skal formaae
at gjøre Rede for, hvad der bliver af Sjælen, naar den forlader
Legemet. Der er i Videnskabens Navn ført strenge Beviser for,
at Sjælen er udødelig, og der er til Gjengjæld ført lige saa strenge
Beviser for, at det Sjælelige selv maa opløses og forsvinde ved
Legemets Død. Begge Forsøg ere lige ugyldige; Videnskaben vil
i begge Tilfælde overskride sin egenlige Grændse. Paa tilsvarende
Maade forholder det sig med den videnskabelige Behandling af
Spørgsmaalet om Guds Tilværelse. Vel kunde det synes, som om
den i Videnskabslæren indførte Subjectivitet, Sandhedens Eenhed
af Viden og Magt, der afgiver den ideel-reelle Forudsætning for
al objectiv virkelig Natur, maatte videnskabelig lade sig bestemme
som Gud selv, som Verdens frie Ophav, som absolut Personlighed
i oprindelig Modsætning til det Upersonlige; men at en Charak\-%
teristik af Grundforholdet mellem Subjectivt og Objectivt, der gaaer
igjennem al Tilværelse, ingenlunde kan aabne Videnskaben noget
Indblik i Guds oververdslige Væsen, i Gud som Gud med absolut
Personlighed, er udtrykkelig fremhævet. „Vi kunne ikke tale om
Aand og Natur, som om Aanden var udenfor Naturen, eller Naturen
udenfor Aanden. Den Eenhed, hvorom her handles, er en Selv\-%
eenhed, der væsenlig hører Naturen til, og uden hvilken Naturen
% doubtful: "saa-vidt" -- a short raised mark (probably a weakly inked hyphen) followed by a space; rejected: "saa vidt"
umulig kan være. For saa-vidt er da ogsaa Selvviden Naturviden,
Selvbegriben Naturbegriben, Selvmagt Naturmagt; men, vel at
mærke, ikke fra Realitetens, men fra Idealitetens Synspunkt.
Vistnok ere Idealitet og Realitet to Sider af det Samme, men disse
to Sider ere ogsaa to Totaliteter i een Totalitet. Den ideelle
Totalitet er Naturens Aand, den reelle Totalitet er den virkelige
Natur. Om man vil kalde denne Anskuelse Theisme, Naturalisme,
Dualisme eller Monisme, er ganske ligegyldigt. Det gjælder om
nøiagtige Bestemmelser, der kunne og maae optages i Aandsvi\-%
% --- p. 43 ---
\opage{43}denskabens Grundlag. Spørgsmaalet, hvorpaa Alt kommer an,
bliver altsaa dette, hvor vidt kan Videnskaben naae, ved Bestem\-%
melsen af Grundforholdet mellem Begreb og Realitet, Selvviden
og Viden af Andet, Selvmagt og Magt i Andet, eller, anderledes
udtrykt, mellem Frihed og Nødvendighed, Aand og Natur“\footnote{Videnskabslære S.\ 227--228.}. % *)

Er det nu Videnskaben lige saa umuligt at bevise Guds Til\-%
værelse som Ikke-Tilværelse, lige saa umuligt at bevise Realiteten
af saadanne personlige Egenskaber som Retfærdighed, Kjærlighed,
Barmhjertighed i den evige Subjectivitet, som at modbevise den,
lige saa umuligt at bevise som at modbevise det Overnaturlige,
skjøndt en Løsning af slige Problemer har afgjørende Betydning
for Menneskets personlige Selv: hvad følger saa heraf? Heraf
følger ganske vist, at der foruden den videnskabelige Erkjendelse,
betinget af Nødvendighed, maa være en religiøs Erkjendelse, be\-%
tinget af Frihed. Det gaaer her med Videnskaben, som det for\-%
hen gik med Religionen. Da man i Religionens Navn havde
gjort Troesprincipet ene gjældende og den kirkelige Dogmatik til
høieste Videnskab, vendte Bladet sig: Fornuften tog Magten fra
Troen, og Erkjendelsesprincipet blev videnskabeligt. Men naar
man saa i Fornuftens Navn vil gjøre Vidensprincipet til enegjæl\-%
dende, altomfattende Princip, saa vender Bladet sig igjen, og vi
indsee, at der ogsaa i Troen maa være et Erkjendelsesprincip, at
der altsaa maa være Noget, som er Gjenstand for Tro, og umulig
kan blive Gjenstand for Viden.

% print: a short centred rule separates Section III from the head; head numeral in heavier type
\begin{center}IV.\end{center}

De to Principer, Troens og Videns, ere, efter hvad vi nu have
seet, ikke at betragte som to Sider af det Selvsamme, ei heller
% --- p. 44 ---
\opage{44}lade de sig opfatte som Brudstykker af eet Princip, der uheldig\-%
viis er revnet og gaaet itu. De to Principer have hvert sit Om\-%
raade, Troesprincipet det Overnaturlige, Vidensprincipet det Natur\-%
lige; og dog have de begge, ethvert især paa sin Viis, samme
Omraade, nemlig Alt, hvad der lader sig sandse og forestille, er\-%
fare og tænke. Hele Troens Indhold kan, ifølge vort Foregaaende
gjøres til Gjenstand for Viden, og omvendt det hele Vidensindhold
faae Betydning for Troen, uden at Troen derfor bliver begrundet
paa Viden eller Viden paa Troen. Det hedder, (Psalm.\ CIV, 14);
„Han“ --- Gud --- „lader Græs fremspire til Qvæget og Sæd til
Gavn for Mennesker“; den Troende fæster sit Blik paa det Natur\-%
lige, for deri at see en Aabenbaring af det Overnaturlige; den
Vidende seer i det Overnaturlige en overflødig Indbildning, da
baade Græs og Sæd kunne voxe frem af naturlig Grund, paa na\-%
turlige Betingelser. Den principielle Forskjel mellem Tro og Viden
kan altsaa ikke bestaae deri, at den Troende og den Vidende slet
ikke forstaae hinanden: men Forstaaelsen bøier sig om i Modsigelse,
idet de, hver efter sit Princip, objectivere paa forskjellig Maade:
i Troen har det Personlige og dermed det Subjective afgjørende
Overvægt, i Viden gjælder det derimod om det Objective, det
Upersonlige.

Dette viser sig klart i Opfattelsen af Love og Mirakler.
Naar man paa Grund af den Umiddelbarhed, hvormed Aabenba\-%
ringstroen i sit Sprog stiller Mirakler imod Naturbegivenheder,
har villet paastaae, at en saadan Mirakeltro nødvendig maa for\-%
udsætte, at Lovene brydes, da er dette, hvad Principet angaaer,
aldeles falsk; Sandheden er, at Troen henfører baade Naturlovene
og Miraklerne til Guds almægtige Villie. Den theologisk-viden\-%
skabelige, d.\ v.\ s.\ uvidenskabelige Opfattelse af Guds Almagt har
ganske vist ført til absurde Conseqvenser\footnote{Malebranche. De la recherche de la vérité p.\ 390. Dieu a créé le
% the third point after "monde" is damaged and looks like a raised comma; read as the third of three points
monde . . . et produit ainsi tous les effets, que nous voyons arriver,
parce qu’ il a voulu aussi certaines loix, selon lesquelles les mouvemens
se communiquent à la rencontre des corps: et parce que ces loix sont
efficaces, elles agissent et les corps ne peuvent agir. Om Theologiens
Naturbegreb. Ved Reformationsfesten 1855.}; % *)
men det er jo ikke
% --- p. 45 ---
\opage{45}om Theologien, det er om Religionen vi her tale. Følger nu
Videnskaben conseqvent sit Princip, maa den unægtelig holde sig
strengt til uforanderlige Love; men da disse Love paa een Gang
tilhøre Viden og Magten, og da Magtens Eenhed med Viden i
Sandhed kan forudsætte Villie, er det Videnskaben umuligt at afgjøre,
om ikke maaskee Lovenes Nødvendighed er i sidste Instans bestemt
ved guddommelig Frihed, altsaa om ikke Miraklernes tilsyne\-%
ladende Afbrydelse af Naturlovene netop er det sande Udtryk for
disses Oprindelse. Videnskaben har, ifølge sit Princip og den
derved bestemte Objectivering, Grund til at forkaste alle Mirakler
uden Undtagelse, men, vel at mærke, paa den Betingelse, at den
anerkjender et andet Princip, for saa vidt den anerkjender en
Mulighed, den ikke selv kan gjennemskue. Videnskaben er ikke
atheistisk, for saa vidt den reent ud erklærer, at den ikke kjender
Religionens Gud, thi denne Gud er ingen videnskabelig Størrelse;
Videnskaben er ikke vantro, fordi den ikke vil vide Noget af
Mirakler, thi Mirakler ere ikke Vidensobjecter. Naar Videnskaben
underkaster Beretningerne om Skabelse og Paradiis en skarp
Kritik, da gjør den kun Rede for sin Erkjendelses Grændse, men
beviser derved ikke, at Skabelse og Paradiis ingen virkelig Be\-%
tydning have.

Har nu Troen sit Princip lige saa vel som Viden, saa følger
jo heraf, at der maa gives Troestanker lige saa vel som Videns\-%
tanker, Troesgrunde lige saa vel som Vidensgrunde, selvfølgelig
ogsaa en Troeskritik lige saa vel som en Videnskritik. Troes\-%
kritiken maa da sikkre den religiøse Opfattelse af Aabenbaringens
Indhold, idet den ikke blot afviser Videns ubeføiede Indblanding,
men ogsaa ved Bestemmelse af Forholdet mellem Mirakel og
Naturlov vogter sig for at komme i Modsigelse med sig selv.
At ogsaa Troen maa have sin Kritik, viser sig da især i dens
Forhold til Miraklerne. Vil Nogen paastaae, at en Troende, der
kan antage eet eneste Mirakel, maa kunne antage alle, da maa
dette vistnok indrømmes, for saa vidt Tro er modsat Viden, men
heraf følger ingenlunde, enten at alle mulige Mirakelfortællinger
skulde i Troens Øine staae paa lige Trin, eller Miraklerne i det
% --- p. 46 ---
\opage{46}Nye Testamente være lige afgjørende for Troeslivet i samme Om\-%
fang; kun Skabelsesmiraklet, Incarnations- og Inspirationsmiraklet
ere for Troeskritiken fundamentale.

Spørgsmaalet om Mirakler og Naturlove hænger paa det
Nøieste sammen med Spørgsmaalet om Fornuftvirkelighed og
Aandsvirkelighed, dette igjen med Spørgsmaalet om det Objective
og dets Objectivering. Hvorfor Paradiis og Syndefald umulig kan
hævde sig virkelig Objectivitet i Videnskaben, er indlysende; hvor\-%
for Paradiis og Syndefald nødvendig maa hævde sig virkelig Ob\-%
jectivitet i Troen, er ligeledes indlysende. Ved at tillægge Troen
og Viden hver sit Erkjendelsesprincip synes vi altsaa at gjøre
Modsætningen saa stærk, at den nødvendig maa sprænge den
menneskelige Bevidstheds Eenhed; den Troende kan ikke være
Vidende, den Vidende ikke Troende; saa ender da det Hele i en
utaalelig Dualisme. „For det udenfor staaende Subject“, siger
% sic: the note's opening „ before "Om Holbergs" has no closing mark
Strauss\footnote{Jevnfør „Om Holbergs Kirkehistorie og Theologi af nærv.\ Forfatter.
Kjøbenhavn 1867. S.\ 71 flgd.}, % *)
„maa Aabenbaringen ved ydre objective Tegn bevise,
at den er Aabenbaring; den maa documentere sin Oprindelse fra
den absolute Intelligens derved, at den paa det Gebeet, som den
endelige Aand kan controllere, viser, hvorledes den overgaaer
denne, hvilket skeer ved \emph{Spaadomme}; den maa bevise sin Op\-%
rindelse fra den absolute Magt ved Herredømme over Natur\-%
magten, ved \emph{Undere}. For at føres til Efterverdenen maa Aaben\-%
baringen bevares i \emph{hellige Skrifter}, men disse maae, for at
man kan være sikker paa, at Aabenbaringen gjemmes uforvansket
i dem, indgives af Aabenbaringens Ophav. Dog dette hele Appa\-%
rat til at gjøre os Aabenbaringen sikker og tilgængelig ligger
endnu bestandig paa den objective Side; der mangler endnu det
Vigtigste, Subjectets Tilegnelse af det aabenbarede Indhold. Sub\-%
jectet tilegner sig Aabenbaringen ved \emph{Skriftfortolkningen},
men hvorledes vil det forvisse sig om, at den foregivne Aabenbaring
virkelig er guddommelig? Ved Underne og Spaadommene, som
ledsagede den, da den fandt Sted. Godt, men hvorledes vil man
% --- p. 47 ---
\opage{47}kunne vide, at disse ere virkelige Undere og Spaadomme? Af
Skriftens Vidnesbyrd. Men hvoraf veed man, at det er sandt?
Deraf, at det er indgivet af den ufeilbare Gud. Men at Skriften
er indgivet af Gud, hvoraf veed man det? Af den \emph{hellige
Aands indre Vidnesbyrd}, som, naar vi læse Skriften, i den
gjenkjender sit eget Værk. Godt, men hvad skal nu overbevise
os om, at dette Vidnesbyrd i os virkelig hidrører fra den hellige
Aand, og ikke fra vor egen, eller endog fra en forførende Aand
udenfor os? --- Her brister det orthodoxe Systems Traad... Den,
som ikke har Aanden i sig, har den udenfor sig; den, som ikke i
sig selv kan finde det bestemmende Princip, søger det i en
Aabenbaring; den, som endnu ikke er moden til Fornufttro, han
bliver ved Aabenbaringstroen“. Stolende paa, hvad han kalder
Fornufttroen, udpeger Strauss en Kløft mellem to Klasser af det
menneskelige Selskab de Vidende og Folket, de Ikke-Philosophe\-%
rende saavel i de høiere som i de lavere Stænder, en Kløft, der
maaskee aldrig vil fyldes. „Saa lade“ siger han, „da den Troende,
den Vidende, ligesom denne hiin roligen vandre sin Vei; vi lade
dem beholde deres Tro, saa maae de ogsaa lade os beholde vor
Philosophi; og hvis det skulde lykkes de overvættes Fromme at
udelukke os fra deres Kirke, saa ville vi agte det for en Vinding;
af falske Foreningsforsøg er der nu skeet nok; kun Skjærpelse af
Modsætningerne kan føre videre.“

At den Straussiske Adskillelse mellem Troende og Vidende i
Bund og Grund er falsk, kan sees af vort Foregaaende. Fornuften
har, naar vi holde os til Principet, hverken nogen Religion ei
heller Erstatning for Religion at byde os; enhver religiøs Gjen\-%
stand, selv en ægyptisk Apis, forudsætter en troende Subjectivitet
og peger i al sin Naturlighed hen paa det Overnaturlige. Har
dette Overnaturlige intet høiere Subjectivt eller Personligt i sig,
maa det conseqvent fortære Menneskets personlige Selv, og Troen
viser sig da virksom fra negativ Side. Den gjennemførte Nihi\-%
lisme, der opløser al Tro, opløser ogsaa al Viden. Dette sees
klart af „Det Ubevidstes Philosophi“\footnote{Philosophie des Unbewussten. Versuch einer Weltanschauung von
Ev.\ Hartmann, Dr.\ phil.\ Berlin 1869.}. % *)
Lige fra Schellings første
% --- p. 48 ---
\opage{48}ungdommelige Optræden med den intellectuelle Anskuelse af det
Absolute indtil den seneste Tid har den speculative Philosophi i
tvetydig Forstaaelse med Naturlæren arbeidet paa at klare sig
det absolute Væsen. Det har imidlertid været det intelligente,
medphilosopherende Publikum vanskeligt at følge med og komme
efter, hvordan det egenlig hang sammen med dette Absolute.
Snart blev det fremstillet som en logisk alvidende Guddom: Gud
i Logiken, før han skabte Verden, snart igjen som et uoplukket,
dunkelt Princip, Altilværelsens iboende Grund og Aarsag. Nu er
da endelig Tvetydigheden hævet; Sandhedsgudens Aasyn sees fuld\-%
stændig afsløret. Verdensgaaden er, at en ubevidst Villen fra
Begyndelsen af har drevet en ligeledes ubevidst Forestilling fra
% sic: the quotation opened at „i sig havde has no closing mark of its own (4 „ against 3 “ on this page)
Ikke-Væren ind i Væren, thi hedder det S.\ 607, „i sig havde
Forestillingen ingen Drift og ingen Interesse af at træde fra Ikke-%
Væren over i Væren, følgelig var før Indtræden af Villen ingen
Forestillen aktuel, følgelig var der før Verdens Tilblivelse hverken
Villen eller Forestillen d.\ e.\ „gar Nichts“. Saalænge Villen ved\-%
varer, saalænge vil Processen og dens Tilsynekommen i Bevidst\-%
heden, Verden, vedvare; naar der altsaa en Gang ikke skal være
nogen Verden mere, da tør der heller ikke mere være nogen
Villen eller nogen Forestilling, da den ubevidste Forestilling altid
kun for saa vidt bliver aktuel, som Villens Interesse fordrer den
d.\ e.\ der vil igjen være „gar Nichts“. Hermed stemmer da ogsaa
følgende Beskrivelse af Guddommen (S.\ 606): „Det Al-Eneste
Ubevidste er alvidende og alviist, og kan følgelig ikke blive klo\-%
gere end det allerede er; det har, som ogsaa Aristoteles siger,
ingen Hukommelse, altsaa kan det af de Erfaringer, det ellers
gjorde i Verden, Intet lære. Følgelig er det, naar Verden en
Gang har hørt op at være, akkurat forblevet det samme, som det
var før Verdens Skabelse; saa saligt som det forhen var, er det
nu igjen, hverken mere eller mindre; Verdensprocessen kan paa
ingen Maade forhjælpe det til større Salighed, end det forhen
havde“. Er Guddommen salig i Intet, da følger det af sig selv,
at Sjælen ogsaa maa blive salig i Intet. Haabet om en Sjælens
individuelle Vedvaren efter Døden erklæres ligefrem for en Illusion.
% --- p. 49 ---
\opage{49}„Dermed er de christelige Forjættelsers Hovednerve overskaaret,
thi Mennesket er det i Grunden kun om hans eget kjære Jeg at
gjøre; hvad hjælper mig den største tilkommende Salighed, naar
jeg ikke kan fornemme og nyde den!“\footnote{For Ide og Virkelighed. Aarg.\ 1875. Hvorledes man danner sig en
Verdensanskuelse. S.\ 12--14.} % *)

I Nihilismen gaaer ikke blot Troen, men ogsaa Viden til
Grunde. At der af Intet selv, det rene Intet, skulde have dannet
sig en Forestillen og Villen og ved deres Forening en virkelig
Verden, som igjen vil opløse sig i det uendelige Intet, er vel saa
meningsløst som muligt. En virkelig Tilværelse, en oprindelig
Modsætningseenhed af Subjectivt og Objectivt og for Menneskets
Vedkommende en væsenlig Forskjel mellem upersonlig Objectiveren
i Viden og personlig Objectivering i Tro, maa nødvendig forud\-%
sættes. Vi skulle ved nogle Exempler forsøge at oplyse dette.
At den bibelske Fremstilling af Verdens Skabelse ikke stemmer
med den videnskabelige Opfattelse af Tingenes Verden, indeholder,
naar vi see nærmere til, ingen Modsigelse. Naar eet og samme
Object objectiveres af Tvende under lige Forhold med samme
Sands, f.\ Ex.\ Synet, og dog fremstiller sig forskjelligt, for den Ene
% sic: "om blaat" as printed (for "som blaat")
om blaat, for den Anden som graat, da maa der naturligviis være
en Feil i Synsorganet enten hos den Ene eller den Anden, om
ikke hos dem begge; men naar den Ene blot kan dømme efter
Synet, den Anden derimod efter sin egen Følelse, saa kan man
ikke forlange fuld Overeensstemmelse, selv om Gjenstanden i og
for sig antages at være den samme. Videnskaben, der dvæler
ved naturlige Elementer, Atomer, Urtaage, physiske Aarsager, Be\-%
tingelser og Love, naaer ikke Aanden; dens upersonlige Objecti\-%
vering, hvor fornuftig den end er, maa blive staaende ved det
udvortes Virkelige. I Religionen derimod er Objectiveringen op\-%
rindelig bestemt ved Aandens eiendommelige Medvirkning. Jorden,
Solen, Maanen opfattes vel, som de umiddelbart vise sig for Sands
og Forstand, men det Afgjørende ligger i den Forvisning, som
Troens Aand vækker om deres Oprindelse og deres Forhold til
% --- p. 50 ---
\opage{50}det Subjective, det Personlige. Det maa udtrykkelig fremhæves,
at den videnskabelige Fremstilling af Universet umulig kan komme
i Strid med den religiøse, med mindre Videnskaben gaaer ud over
sin Grændse, afskaffer Aabenbaringsphantasien og selv paatager
sig at afgjøre, hvorledes det principielt forholder sig med Skaber
og Skabning. Heller ikke kan den religiøse Opfattelse af en Ska\-%
belse i sex Dage komme i Strid med Læren om en successiv Ud\-%
vikling, med mindre Religionen vil opgive sin principielle Be\-%
grændsning, sætte sig i Videnskabens Høisæde og klarlig vise,
hvorledes Millioner af Aar svinde ind til sex Dage.

„Men“, indvender man, „skal et Forlig mellem Religion og
Videnskab begrundes paa en saa principiel Forskjel i Objectiverings\-%
maaden, at Videnskaben giver den virkelige, fra alle uvedkommende
Indblandinger befriede Objectivitet, som den er, medens Religionen
faaer sin eiendommelige Verden ved en Objectivering, der omdanner
Virkeligheden, idet den indblander det Personlige, saa er det Hele for\-%
virret. Hvad bliver der vel af dette Personlige, naar vi tale reent ud?
Det Mythiske! Ere nu Udtryk som personlig Objectivitet, Aandsvirke\-%
lighed, Aabenbaringsvirkelighed intet Andet end Omskrivninger af,
hvad vi allerede kjende under Navn af det Mythiske, saa er herved jo
slet Intet vundet til virkelig Belysning af Religionens Selvgyldig\-%
hed over for Videnskaben“. Vi ville prøve denne Indvending ved
at kaste et Blik paa Paradiis. At Paradiset under enhver Form
maa i Videnskaben være en illusorisk Virkelighed, et Phantasiens
Blændværk, en Mythe, kan ifølge vort Foregaaende ikke modsiges;
men vi have jo indrømmet Troen sit Princip, at den ikke
skulde staae under Videns Formynderskab, men have Ret til at
hævde sit eget Syn paa Verden og anvende sin egen Kritik. Lad
os da med Troesprincipet for Øie betragte Paradiset fra Aaben\-%
baringens, Mythens og Poesiens Standpunkt. Aabenbaringen be\-%
gynder med en Omgivelse, hvori Guddommen i sandselig-over\-%
sandselig Nærværelse lever sammen med Mennesket; denne para\-%
disiske Virkelighed er Aandsvirkelighed. Den umiddelbare Sam\-%
mensmeltning af Naturligt og Overnaturligt hører imidlertid alle\-%
rede Forbigangenheden til, naar Paradiset bliver Gjenstand for
% --- p. 51 ---
\opage{51}Tro. Da det Forbigangne ikke kan sees umiddelbart, men maa
anskues i Phantasiens Reflexionsspeil, bliver det uvist, hvorvidt
Troen kan finde de givne Tegn saa vel fra guddommelig som fra men\-%
neskelig Side saa fyldestgjørende, at den i de paradisiske Phæno\-%
mener maa fastholde væsenlige Træk af en sand Aabenbaring,
eller der, idet Reflexionen kommer i Bevægelse, indtræder en
Uvished om, hvad der er guddommeligt og hvad der er menne\-%
skeligt, hvor meget der er overnaturligt, hvor meget der er natur\-%
ligt. I sidste Tilfælde forvandles Aabenbaringen til Mythe. Det mythi\-%
ske Væsen har mange vexlende Skikkelser, og afgiver Grundlag for
Poesi og Kunst. Det Paradiis, hvormed Ovid begynder sine Meta\-%
morphoser, er aabenbart blot poetisk; Paradiset i den persiske Religion
er mythisk; men hvorfor er Hebræernes og de Christnes Paradiis
hverken blot poetisk ei heller mythisk, men en virkelig Aabenba\-%
ring? Det maa principielt afgjøres i Tro, ikke i Viden. Ifølge
Vidensprincipet ere alle Paradiser mythiske, det mosaisk christelige
lige saa vel som det persiske og muhammedanske. Den paradisiske
Troesvirkelighed er Aandsvirkelighed; den Tro, som i dybeste For\-%
stand erkjender det personlige Forhold mellem Gud og Menneske,
maa altsaa i Aabenbaringens Navn erkjende det sande, aands\-%
% a small low point stands between "virkelige" and "Paradiis" (smaller than a full stop; spacing material?); not transcribed
virkelige Paradiis.

Skulle vi særlig undersøge Christi og Djævelens Virksomhed,
da gjælder det samme Princip, Troesprincipet. At der i Palæstina
virkelig har levet et Menneske ved Navn Jesus, som har sagt sig
at være Christus, Guds eenbaarne Søn, vil Videnskaben ikke be\-%
nægte; Jesus af Nazareth er virkelig en historisk Personlighed,
og har som saadan Krav paa historisk Tro. Men historisk Tro
(fides historica) er principielt forskjellig fra religiøs Tro (fides
religiosa), Jesus som Guds menneskeblevne Søn, Jesus som den
Jomfrufødte, som Undergjøreren, som den Opstandne og Himmel\-%
farne, kort: Jesus som Christus er ikke Gjenstand for blot historisk,
men for religiøs Tro. Hvad der gjælder om Christus, gjælder
ogsaa om Djævelen. Vel er Djævelen selv endnu ikke incarneret
--- Antichristen er jo endnu ikke kommen --- dog sees det af vort
Foregaaende, at den videnskabelige Kritik, der begynder med at
% --- p. 52 ---
\opage{52}opløse Djævelens Personlighed, conseqvent maa ende med at op\-%
løse Christi Personlighed; hvoraf da ogsaa det Omvendte følger,
at den Tro, som principielt vil fastholde den personlige Christus,
ogsaa maa erkjende den personlige Djævel.

Vil man nu heraf slutte, at naar Tro og Viden ere saa dia\-%
metralt modsatte, saa kan den Troende umulig være vidende, og
den Vidende ikke troende: da er Slutningen grundfalsk. Ere Troes\-%
objecterne saa principielt forskjellige fra Vidensobjecterne, at En
kan være troende uden at være vidende, fordi de religiøse Objecter
ere personlige, de videnskabelige derimod upersonlige: saa følger
ogsaa heraf, at den Troende kan være en Vidende, og den Vidende
en Troende uden at komme i Modsigelse med sig selv. En Astro\-%
nom, der paa een Gang vil hylde baade det ptolemæiske og det
copernicanske Verdenssystem, kommer ganske vist i Modsigelse
med sig selv, thi begge Systemer ere Resultater af samme Art
Objectivering: hvis det ene er sandt, maa det andet nødvendig
være usandt. Den derimod, der i Egenskab af Astronom hylder
det copernicanske System, men i Egenskab af Troende derimod
den mosaiske Fremstilling, kommer ikke i Modsigelse med sig
selv; thi idet han gjennem det Naturlige sigter paa det Overnatur\-%
lige og objectiverer personligt i Phantasi, kan han umulig objectivere
videnskabeligt. Naar En i Troen fastholder, at Mennesket blev skabt
paa den sjette Dag, medens en Anden troer, at det først er blevet
skabt sex Millioner Aar efter Jordens første Dannelse, da befinde
de sig i virkelig Modsigelse med hinanden, thi de bevæge sig
indenfor samme Sphære; men naar en Troende paastaaer, at Men\-%
nesket er skabt i Guds Billede, medens en Zoolog paastaaer, at
Mennesket har udviklet sig af samme Grundform som Aben: da
er heri ingen Modsigelse; thi der tales om høist forskjellige Ting.
„At Gud Jehova dannede Mennesket, Støv af Jorden, og indblæste
i dets Næseboer Livets Aande“ (1ste Moseb.\ II, 7), angaaer ikke
den naturlige Udvikling, men alene den Oprindelse fra Gud, som
kun kan anskueliggjøres i personlig Objectivering; Menneskets Ud\-%
vikling af en naturlig Grundform afgjør Intet om Skabelsen. Hvis
den Troende derimod ikke vil nøies med at fremhæve den gud\-%
% --- p. 53 ---
\opage{53}dommelige Skaberact, paa hvilken Indblæsningen af Livets Aande
er et betegnende Symbol, men indlader sig paa at bestemme Udvik\-%
lingens Gang, krænker han sit Princip og faaer allerede derved
Uret. Vil omvendt den videnskabelige Forsker, der ifølge Videns\-%
principet objectiverer sig den trinvise Naturudvikling, fornægte det
i Anlæget liggende Skabelsesmoment, han slet ikke kjender, over\-%
skrider han sin Grændse og krænker sit Princip. Her er fra
begge Sider et Rethaveri, som i lige Grad skader Religion og
Videnskab.

Indvender man maaskee, at det kun er i visse særegne Til\-%
fælde, at man kan beraabe sig paa en principielt forskjellig Objec\-%
tiveringsmaade og en tilsvarende forskjelligartet Virkelighed, for
at hæve Striden mellem Tro og Viden, men at denne Udvei er
langt fra at slaae til i de vanskeligste, vigtigste Tilfælde, saa lad
os endnu en Gang kaste et Blik paa Christi Opstandelse og
Himmelfart. Er det virkelig Mennesket Jesus, og ikke en almin\-%
delig Christusidee eller blot aandelig Christus, der er opstanden
fra de Døde og opfaren til Himlen, saa kan Troesobjectet umulig
løsrives fra Vidensobjectet, uden at den hele Virkelighed maa blive
mythisk. Den, der opstaaer af Graven og viser sig for Disciplene,
er jo den Korsfæstede, den, hvem Alle, baade Vantroe og Troende,
kunde see med sandselige Øine og gribe med sandselige Hænder,
altsaa objectivere sig baade paa Videns og Troens Maade; men
efter Opstandelsen er hans Skikkelse som ganske forvandlet, virke\-%
lig og dog overvirkelig; Maria Magdalene kan ikke kjende ham
% sic: "(Luk. XXIV;" has no closing parenthesis
(Joh.\ XX.\ 15--17), de To, der gik til Emaus, heller ikke (Luk.\
XXIV; han gaaer igjennem lukkede Døre, og Disciplene mene at
see en Aand; men Thomas, der vil have haandgribelig Vished,
faaer den og troer (Joh.\ XX; 19--31). Om videnskabelig Frem\-%
stilling af slige Begivenheder kan her slet ikke være Tale. Ob\-%
jectiviteten er som Religionen selv, Aabenbarelsen af det i Person\-%
ligheden Ubegribelige, Aandsvirkelighed i beskuelig Form, Gjen\-%
stand for Tro, ikke for Viden.

At Tro og Viden, hver med sit Princip, ikke staae i Mod\-%
sigelse med hinanden, maa vel af det Foregaaende være klart;
% --- p. 54 ---
\opage{54}men Tro og Viden ere jo ueensartede Erkjendelser: Spørgsmaalet
bliver da, hvorledes de lade sig forene i den menneskelige Be\-%
vidstheds Eenhed? Hertil maa Følgende svares: Videnskaben,
der er bleven sig sit Princip, sin Ret, sin Frihed fuldt bevidst,
kan i Religionen som guddommelig Aabenbaring ikke see nogen
Fjende, der vil gribe forstyrrende ind i dens Undersøgelser; om
end den videnskabelige Erkjendelse ikke kan uddybe Personlig\-%
heden selv, maa der dog være Fornuft i al personlig Tænkning.
Fra sin Side kan Religionen da heller ikke staae fjendtlig imod
Videnskaben; om den end, af personlige Grunde, henfører Til\-%
værelsens Love under Guds Villie og lader det Naturlige være
det Overnaturlige underordnet, saa kan denne Underordnen dog
umulig betyde noget Brud paa Lovene. Naar det hedder, at
Christus har været under Loven, for at fuldkomme Loven, da
gjælder det ikke om Moseloven alene, men om Verdensloven,
% sic: "Skiøndt" as printed
Naturloven. Skiøndt Religionen med sit Personlighedsprincip ikke
er afhængig af Videnskaben, er den derfor ikke udelukket fra
Videnskaben. Religionens Princip er ogsaa Sædelighedens Princip;
men Sædelighedsprincipet virker i Samfundet paa saa naturlige
Betingelser, at det i sin Almeengyldighed maa tilfredsstille For\-%
nuftens Fordringer, og desaarsag er blevet forvexlet med For\-%
nuftprincipet selv.

Ogsaa Religionen eller rettere Religionslæren har Trang til
videnskabelig Belysning, kun at denne Belysning ikke giver sig
Skin af videnskabelig Begrundelse. „Religionsphilosophien bliver
ikke staaende ved almindelige Grændsebestemmelser for Tro og
Viden, den kan umulig lade Religionen udvikle sit Indhold,
% sic: "ude" as printed (for "uden")
ude at lade et heelt System af Troesbegreber komme til Ud\-%
vikling. Men Troens religiøse Indhold er, ifølge Principet,
et overvidenskabeligt Indhold; her indtræder da den Mod\-%
sigelse, at et overvidenskabeligt Indhold optages i Videnskaben.
Dersom Videnskabsformen nu virkelig blev forvexlet med Indholdets
sande, oprindelige Form og indført og fastholdt istedenfor denne,
vilde Modsigelsen ganske vist være uopløselig. Men dette er
ingenlunde Tilfældet. Den videnskabelige Reflexion har just den
% --- p. 55 ---
\opage{55}Bøielighed, at den bringer Vidensformen til at reflectere en anden
og modsat Form. Thi ligesom den klare Dag kan bringe Øiet
til at glemme Lysskjæret og see paa Gjenstandene, saaledes kan ogsaa
Reflexionen bringe Tanken til i Viden at see paa Troen og glemme
Viden. Religionsphilosophien er saa langt fra at paatvinge det
religiøse Indhold en fremmed Form, at dens Formudvikling tvært\-%
imod bestaaer i at frigjøre det fra Tilsætninger af umiddelbart
% sic: "o, s. v." (comma after "o") as printed
individuelle Forestillinger, Billeder, Følelser, Villiesytringer o, s.\ v.,
hvormed det i Troeslivet selv altid er sammenvoxet. For saa
vidt Reflexionen, saavel hvad Formen som hvad Indholdet angaaer,
ved Viden ophæver Viden, idet den lægger an paa at klare, ikke
en videnskabelig, men en overvidenskabelig Sandhed, er Forholdet
mellem Sandheden og Videnskaben væsenlig vendt om, og Reli\-%
gionsphilosophien bliver med Hensyn hertil at charakterisere som
en omvendt Videnskab“.\footnote{Religionsphil.\ S.\ 8--11.} % *)

\emph{I skulle erkjende Sandheden}, og \emph{Sandheden skal
frigjøre Eder}.
% end of Nielsen's text: a short double rule is printed below the footnote on p. 55; no signature or date

\end{document}
